% Open-vocabulary object lists from SAM 3 + VGGT-SLAM (task1-sam3): findings.
% Fragment for % Open-vocabulary object lists from SAM 3 + VGGT-SLAM (task1-sam3): findings.
% Fragment for % Open-vocabulary object lists from SAM 3 + VGGT-SLAM (task1-sam3): findings.
% Fragment for % Open-vocabulary object lists from SAM 3 + VGGT-SLAM (task1-sam3): findings.
% Fragment for \input{sam3}; needs \usepackage{booktabs}.
% All numbers: VSI-Bench, 5,130 questions, 288 scenes (287 processed); box score threshold 0.7 unless stated.
% Sources: VGGT-SLAM/exps/vsibench_objects (s3_alt_stats.py, s3_badlist_thr.py, qvocab.py, s3_rows.py),
% 3d-scene-lang/runs/eval_q35vs_gandalf/{mixv2_1_sam3_full,mixv2_1_s3q_roomdoor}_rvs_question.

\section{Open-vocabulary object lists with SAM 3}
\label{sec:sam3}

\paragraph{Pipeline.}
SAM 3 is run on up to 100 VGGT-SLAM keyframes per video with one text prompt per object name
(per-frame confidence $\geq 0.5$). Each mask is lifted to 3D with the SLAM points, masks of the same name are
merged across frames, and each merged instance gives one box. The box score is the mean of the three best
SAM 3 scores times $1-\exp(-\mathrm{views}/2)$. No ground truth is used. The resulting list of
\{label, box, score\} entries is given to the language model together with the video.

\paragraph{The object list, not the geometry, limits the benefit.}
Replacing the predicted centres and sizes by ground truth while keeping the SAM 3 list raises the VSI-Bench
score only from 68.6 to 69.9, whereas the full ground-truth list reaches 77.3. The errors that matter are
extra boxes and missing objects. With question-name prompts and ``room door'' (below), the list has 414 extra
and 641 missing boxes against 3,194 ground-truth objects of the names the questions use.

\subsection{Run time}
\label{sec:sam3-time}

Table~\ref{tab:sam3-time} gives the cost of the two stages, measured on 597 ScanNet++ scenes processed together
on four NVIDIA RTX PRO 6000 Blackwell Max-Q Workstation Edition GPUs (96\,GB each, machine gandalf) with three
workers per GPU; the GPUs ran nothing else. VGGT-SLAM receives up to 600 frames per scene and
keeps 100 keyframes; SAM 3 runs all 65 class prompts on each of the 100 keyframes (per-frame confidence 0.3).
SAM 3 costs about eight times as much as VGGT-SLAM per scene. Its time hardly varies between scenes, since every
prompt is run on every keyframe (about 4\,s per keyframe for 65 prompts); with three workers sharing a GPU it is
limited by GPU compute, so more workers per GPU do not raise the throughput. For comparison, SAM 3 on ARKitScenes on
four NVIDIA RTX 6000 Ada GPUs (46\,GB each, machine wukong; two workers per GPU) took about
250\,s per scene per worker with 51 prompts.

\begin{table}[t]
\centering
\caption{Run time per scene of the two stages (597 ScanNet++ scenes, 4 NVIDIA RTX PRO 6000 Blackwell Max-Q GPUs,
96\,GB, 12 workers).}
\label{tab:sam3-time}
\begin{tabular}{lrrrr}
\toprule
Stage & \multicolumn{3}{c}{Time per scene per worker} & Throughput \\
      & Median & Mean & 90th percentile & (12 workers) \\
\midrule
VGGT-SLAM (600 frames in, 100 keyframes)   &  48\,s &  48\,s &  74\,s & 10 scenes per minute \\
SAM 3 (65 prompts on 100 keyframes)        & 394\,s & 392\,s & 398\,s & 110 scenes per hour \\
\bottomrule
\end{tabular}
\end{table}

\subsection{Prompt wording changes the number of boxes}

\paragraph{``door'' versus ``room door''.}
The prompt ``door'' also fires on cabinet, wardrobe and shower doors. It alone produced 37\% of all extra boxes
(231 of 620). Prompting ``room door'' instead removes most of them at the cost of some recall
(Table~\ref{tab:sam3-door}).

\begin{table}[t]
\centering
\caption{Door boxes for the two prompts (170 ground-truth doors in 114 scenes).
Found: a box of that label within 0.75\,m of the ground-truth centre.}
\label{tab:sam3-door}
\begin{tabular}{lrrrrr}
\toprule
Prompt & Boxes & Extra & Missing & Scenes with exact count & Found \\
\midrule
``door''      & 391 & 231 & 10 & 31\% & 85\% \\
``room door'' & 163 &  25 & 32 & 64\% & 76\% \\
\bottomrule
\end{tabular}
\end{table}

\paragraph{A more specific word is not always better.}
For tables, every rewording we tried is worse than the plain name (Table~\ref{tab:sam3-table}): ``desk'' removes
13 extra boxes but loses 61 more tables. The extra table boxes are duplicates of real tables and tables without
annotation, so they call for better merging, not for another word.

\begin{table}[t]
\centering
\caption{Table boxes for four prompts (406 ground-truth tables, 285 scenes, all with ``room door'').}
\label{tab:sam3-table}
\begin{tabular}{lrrrrr}
\toprule
Prompt & Boxes & Extra & Missing & Scenes with exact count & Found \\
\midrule
``table'' (benchmark name)      & 450 & 101 &  57 & 41\% & 75\% \\
``desk''                        & 376 &  88 & 118 & 35\% & 62\% \\
``table or desk''               & 529 & 162 &  39 & 34\% & 79\% \\
``dining table or office desk'' & 287 &  58 & 177 & 36\% & 50\% \\
\bottomrule
\end{tabular}
\end{table}

\paragraph{Which names are detected too rarely or too often.}
Table~\ref{tab:sam3-names} lists the names with the largest count errors. Two groups appear.
Names detected too rarely are mostly ones whose benchmark wording is unusual for the object
(``heater'' for a wall radiator, ``ceiling light'', ``computer tower'') or small objects (basket, shoes, bucket,
pan, kettle), which have no extra boxes at all. Names detected too often (table, chair, sofa, bed) are large
furniture that is split into several boxes or duplicated. ``pillow'' and ``cushion'' fire on the same objects.
Only the door and table wordings have been tested; the others are candidates.

\begin{table}[t]
\centering
\caption{Count errors per object name (question-name prompts with ``room door''). Scenes: those whose questions
use the name and that contain the object.}
\label{tab:sam3-names}
\begin{tabular}{lrrrrr}
\toprule
Name & Ground truth & Boxes & Extra & Missing & Scenes with too few \\
\midrule
\multicolumn{6}{l}{\emph{Too few boxes}} \\
heater / radiator & 46 & 6 & 0 & 40 & 82\% \\
ceiling light     & 116 & 53 & 1 & 64 & 72\% \\
computer tower    & 42 & 25 & 1 & 18 & 71\% \\
fan / exhaust fan & 14 & 5 & 1 & 10 & 60\% \\
power strip       & 10 & 6 & 1 & 5 & 57\% \\
backpack          & 43 & 29 & 0 & 14 & 48\% \\
window            & 83 & 68 & 11 & 26 & 40\% \\
\midrule
\multicolumn{6}{l}{\emph{Too many boxes}} \\
table & 409 & 454 & 102 & 57 & 23\% \\
chair & 897 & 864 & 73 & 106 & 26\% \\
sofa  & 146 & 144 & 30 & 32 & 31\% \\
bed   & 82 & 106 & 26 & 2 & 3\% \\
door (``room door'') & 170 & 163 & 25 & 32 & 18\% \\
\bottomrule
\end{tabular}
\end{table}

\paragraph{The score threshold trades missing objects for extra ones.}
Lowering the box-score threshold does not fix the list (Table~\ref{tab:sam3-thr}, fixed 65-name prompts): at 0.3
fewer labels are short of boxes, but there is about one extra box per ground-truth object. About 7\% of labels
are never found at any threshold, because SAM 3 gives no mask above its per-frame confidence.

\begin{table}[t]
\centering
\caption{Per (scene, label) count against ground truth for four score thresholds.}
\label{tab:sam3-thr}
\begin{tabular}{rrrrr}
\toprule
Threshold & Exact & Too few & Too many & Extra boxes per ground-truth object \\
\midrule
0.7 & 44.4\% & 22.6\% & 33.0\% & 0.43 \\
0.5 & 41.6\% & 15.6\% & 42.8\% & 0.64 \\
0.3 & 30.9\% &  9.9\% & 59.2\% & 1.06 \\
0.2 & 24.2\% &  7.0\% & 68.8\% & 1.48 \\
\bottomrule
\end{tabular}
\end{table}

\subsection{The vocabulary of the questions}

\paragraph{Seven of eight tasks use a closed list of 65 names.}
Outside route planning, every object in a VSI-Bench question is one of 65 names, and the wording depends on the
source dataset (``trash can'' / ``trash bin'', ``heater'' / ``radiator'', ``whiteboard'' / ``blackboard'').
Prompting SAM 3 with the names a scene's own questions use, in their original wording, gives 15.9 boxes per
question instead of 41.9 for a fixed list of all names.

\paragraph{Route-planning questions are open-vocabulary.}
The 194 route-planning questions are free text. Of their 979 object mentions (start, facing direction, goal and
every ``go forward until''), only 505 (52\%) are one of the 65 names. The rest are other objects
(sink in 28 scenes, cabinet 11, bathroom sink 8, desk 6, doorframe 6, wardrobe 4),
objects with an attribute (``blue chair'', ``open door'', ``glass coffee table''),
spelling variants (``trashbin'', ``bathtube'', ``refridgerator'', ``dinning table'')
and places that are not objects (wall, ``into the room'', ``outside'').
A fixed vocabulary therefore gives no box at all for about half of the objects these questions refer to.

\paragraph{Extraction rule.}
We take each mention in the question's own words, remove articles, position words and trailing relations
(``the sink on the left'' $\to$ ``sink'', ``foot of the bed'' $\to$ ``bed''), keep attributes
(``blue chair''), and drop walls, rooms and ``outside''. This gives 183 distinct phrases; 241 (scene, phrase)
pairs in 147 scenes are not already question names and become additional prompts. Only a plain ``door'' is
prompted as ``room door''; ``doorframe'', ``doorway'', ``open door'' or ``kitchen door'' are used as written.
When a route phrase and a benchmark name find the same object, the list keeps one box, shown under the name the
question uses.

\subsection{Effect on VSI-Bench}

Table~\ref{tab:sam3-vsi} compares the lists with the same language model (Qwen3.5-4B, mix-v2-1).
Question-name prompts with ``room door'' gain 0.4 / 0.5 points over the fixed 65-name prompts. The gain is in
counting ($+1.5$) and absolute distance ($+0.8$), the tasks most affected by extra door boxes; the other tasks
change by at most 0.6. The list still does not beat the video alone.
Adding the route phrases changes the boxes shown for 130 of the 194 route-planning questions, yet the
route-planning score stays at 52.1: only 10 answers change, 4 from wrong to right and 4 from right to wrong.
Section~\ref{sec:sam3-route} examines why.

\begin{table}[t]
\centering
\caption{VSI-Bench score (mean over tasks) on all 5,130 questions and on the debiased subset.}
\label{tab:sam3-vsi}
\begin{tabular}{lrrrrr}
\toprule
Object list & Full & Debiased & Counting & Abs.\ distance & Route planning \\
\midrule
None (video only)                               & 69.4 & 63.7 & -- & -- & -- \\
SAM 3, fixed 65 names                           & 68.6 & 62.3 & 72.0 & 56.6 & 51.5 \\
SAM 3, question names + ``room door''           & 69.0 & 62.8 & 73.5 & 57.4 & 52.1 \\
SAM 3, question names + ``room door'' + route phrases & 69.0 & 62.7 & 73.3 & 57.4 & 52.1 \\
SAM 3 list with ground-truth centre and size    & 69.9 & 63.9 & -- & -- & -- \\
Ground-truth list                               & 77.3 & 73.9 & -- & -- & -- \\
\bottomrule
\end{tabular}
\end{table}

\subsection{Route planning: are the named objects in the list?}
\label{sec:sam3-route}

Route planning is the task an open vocabulary should help most, since half of its object mentions are outside
the 65 names. We therefore check, for each of the 194 questions, whether every object it names has a box in
the list that includes the route phrases. A question names 659 distinct objects in total, 418 by a benchmark
name and 241 by another phrase.

\begin{table}[t]
\centering
\caption{Route-planning questions: objects named in the question that have at least one box of that name,
for four box-score thresholds. The evaluated list uses 0.7.}
\label{tab:sam3-route}
\begin{tabular}{rrrr}
\toprule
Threshold & Questions with every object found & Benchmark names found & Other phrases found \\
\midrule
0.7  & 120 / 194 (62\%) & 375 / 418 (90\%) & 198 / 241 (82\%) \\
0.5  & 145 / 194 (75\%) & 394 / 418 (94\%) & 210 / 241 (87\%) \\
0.3  & 157 / 194 (81\%) & 400 / 418 (96\%) & 218 / 241 (90\%) \\
none & 164 / 194 (85\%) & 403 / 418 (96\%) & 222 / 241 (92\%) \\
\bottomrule
\end{tabular}
\end{table}

\paragraph{Coverage is incomplete.}
At the evaluated threshold only 62\% of the questions have a box for every object they name
(Table~\ref{tab:sam3-route}). The objects missing most often are heater (9 questions), door (7), window (4),
cabinet (3) and shower (3). Heater is also the main object never found at any threshold (8 questions),
in line with Table~\ref{tab:sam3-names}; the remaining never-found objects are single cases such as
``column'', ``trophies'' or ``paper shredder box''.

\paragraph{Found objects are often ambiguous.}
A route question refers to one specific object (``the door'', ``the blue chair''). At threshold 0.7,
149 of the 375 found benchmark names and 91 of the 198 found phrases have several boxes of that name, so the
list does not say which box is meant. Lowering the threshold closes part of the coverage gap but widens this
one: with no threshold, 295 of the 403 found benchmark names have several boxes.

\paragraph{The language model barely uses the added boxes.}
The route phrases change the boxes shown for 130 of the 194 questions, and the 6 questions that matched no box
before now match some. Yet only 10 answers change, 4 from wrong to right and 4 from right to wrong; 101 answers
are correct with and without the phrases (52.1\%). A likely reason, which we have not
tested, is that the model was fine-tuned on lists that use the dataset label names, so boxes under free phrases
are an input it has not been trained to use.

\paragraph{Limits of this check.}
``Found'' means that a box with the name exists. Free phrases have no ground-truth boxes, so we cannot verify
that the box is on the right object. Mentions that are not objects (walls, rooms, ``outside'') are not counted.

\paragraph{Conclusion for route planning.}
Three gaps act together: 38\% of the questions miss at least one named object, about 40\% of the found objects
are ambiguous, and the model does not exploit the additional boxes. Extending the vocabulary alone therefore
gives no gain (52.1 with and without route phrases); the task needs boxes that identify the referred instance
and a model trained on such lists.

\paragraph{Summary.}
(i) Text prompts are not interchangeable: one word (``room'') removes 206 of 231 extra door boxes, while a
plausible synonym for table loses 61 objects.
(ii) Benchmark wording should be kept per dataset, and ambiguous names disambiguated, not replaced by synonyms.
(iii) A closed vocabulary covers only about half of the objects in free-text questions; taking the prompts from
the question raises coverage to 82--90\% of the named objects, but does not change the route-planning score
(Section~\ref{sec:sam3-route}).
(iv) Wording fixes a minority of the errors: after them, missing objects (641) outnumber extra boxes (414), and
neither a lower threshold nor rewording removes both.

\subsection{Open questions}
\label{sec:sam3-open}

\paragraph{Q1. Do the 65 classes (with ``room door'') cover the detector's training data?}
The detector (VGG-T-Det) is trained on ScanNet, ScanNet++ and ARKitScenes. Do the 65 class names, prompted as
in Section~\ref{sec:sam3} with ``room door'' for door, cover the annotated objects of these three datasets, and
which annotated objects fall outside them?
\emph{Known so far:} the 65 classes are the VSI-Bench names, so every class has a benchmark use, but the
converse is not measured. In the 242 ScanNet++ scenes that are ready, 7,881 of 27,878 annotations map to one of
the 65 classes (the rest include walls, floors and the generic ``object''); seven classes have boxes only in
ScanNet++ (ceiling light, cutting board, paper bag, shoe rack, coat rack, power strip, mattress) and two have
none there (fireplace, piano). The shares for ScanNet and ARKitScenes, and whether SAM 3 with these prompts
finds the annotated objects of each class in the training scenes, are still to be measured.

\paragraph{Q2. Do the 65 classes cover the objects named in the Reason-190K questions?}
The language model is fine-tuned on Reason-190K. If its questions name objects outside the 65 classes, a
65-class list gives no box for them in training, as it does for route planning in VSI-Bench.
\emph{Known so far:} not measured for the 65-class set. The measurement is the one of
Section~\ref{sec:sam3-route} applied to the training questions: the share of questions whose named objects are
all among the 65 names, per source dataset.

\paragraph{Q3. Should the whole pipeline be open-vocabulary?}
If the answers to Q1 and Q2 show large gaps, the alternative is to drop the fixed class set: the label of a box
comes from the SAM 3 prompt, the classification head of the detector is removed, and the detector only refines
the centre and size of each box and decides whether to keep it. For this to work the detector must receive
more candidate boxes than are finally kept, so the SAM 3 score threshold would be lowered and the removal of
extra boxes left to the keep decision.
\emph{Known so far:} lowering the box-score threshold from 0.7 to 0.3 reduces the labels with too few boxes
from 22.6\% to 9.9\% but raises the extra boxes per ground-truth object from 0.43 to 1.06
(Table~\ref{tab:sam3-thr}), so a keep decision would have to remove about one box per object. About 7\% of
labels are found at no box-score threshold; recovering them needs a lower SAM 3 per-frame confidence (0.5 now),
which has not been tried. Two risks are open: a correct box on an object without annotation is trained as
``drop'', and the detector cannot add an object that SAM 3 missed.

\paragraph{Q4. Should SAM 3 run on overlapping video segments that are merged afterwards, like submaps?}
VGGT-SLAM handles long videos by reconstructing submaps of consecutive frames that share overlapping frames and
then aligning them. The same scheme could be used for the objects: run SAM 3 on each video segment, with the
overlapping frames in common, build the object instances per segment, and merge the segment instances at the end,
using the overlapping frames to decide which instances are the same object.
\emph{Known so far:} the current pipeline does not use segments. SAM 3 is applied as an image model to each
keyframe independently (at most 100 keyframes spread evenly over the video, so long videos are subsampled), and
masks are associated only afterwards in 3D, by the overlap of their lifted points over the whole video. Extra
boxes from this global merge are the main count error for large furniture (Table~\ref{tab:sam3-names}), and
errors of the global reconstruction enter every association. Open points: whether SAM 3's video mode, which
tracks an object through consecutive frames, gives cleaner instances within a segment than per-frame masks;
whether merging per segment and then across segments produces fewer duplicates than one global merge; and how
many keyframes per segment are needed so that small objects are seen in several views.

\paragraph{Q5. Which count errors can be fixed by wording, and which need better merging?}
Two groups of names have wrong box counts (question-name prompts with ``room door'', box score $\geq 0.7$).
For the names with too few boxes, the benchmark wording is the likely cause, and
Table~\ref{tab:sam3-q5-few} lists a wording to try for each. For the names with too many boxes, the likely
cause is the merging of masks into instances, not the wording (Table~\ref{tab:sam3-q5-many}).
\emph{Known so far:} only the door and table wordings have been tested
(Tables~\ref{tab:sam3-door} and~\ref{tab:sam3-table}). The wordings and causes in the two tables below are
hypotheses read from the counts; none has been tested. Small objects with too few boxes and no extra boxes
(basket, shoes, bucket, pan, kettle, blanket, bowl, cup) are probably not a wording problem; a lower SAM 3
per-frame confidence is the more likely remedy.

\begin{table}[t]
\centering
\caption{Names with too few boxes, where the wording is the likely cause, and the wording to try (untested).}
\label{tab:sam3-q5-few}
\begin{tabular}{lrrrl}
\toprule
Name & Ground truth & Boxes & Scenes with too few & Wording to try \\
\midrule
heater / radiator & 46  & 6  & 82\% & ``radiator'' for both, or ``wall radiator'' \\
ceiling light     & 116 & 53 & 72\% & ``ceiling lamp'' (the ScanNet++ annotators' word) \\
computer tower    & 42  & 25 & 71\% & ``desktop computer'', ``PC case'' \\
fan / exhaust fan & 14  & 5  & 60\% & ``ventilation fan'' \\
power strip       & 10  & 6  & 57\% & ``extension cord with sockets'' \\
counter           & 13  & 12 & 55\% & ``kitchen counter'' \\
backpack          & 43  & 29 & 48\% & ``bag'' as a second prompt \\
window            & 83  & 68 & 40\% & unclear; may be visibility, not wording \\
\bottomrule
\end{tabular}
\end{table}

\begin{table}[t]
\centering
\caption{Names with too many boxes, where merging rather than wording is the likely cause (untested).}
\label{tab:sam3-q5-many}
\begin{tabular}{lrrrl}
\toprule
Name & Ground truth & Boxes & Extra & Likely cause \\
\midrule
table            & 409 & 454 & 102 & duplicates, tables without annotation \\
chair            & 897 & 864 & 73  & duplicates (also 106 missing) \\
sofa             & 146 & 144 & 30  & sectional sofas split into pieces \\
bed              & 82  & 106 & 26  & bed split into several boxes \\
tv               & 94  & 96  & 15  & monitors detected as tv \\
pillow / cushion & 100 & 105 & 20  & the two names detect the same objects \\
\bottomrule
\end{tabular}
\end{table}

\paragraph{Q6. Which box-score threshold is best for the language model?}
All VSI-Bench results above use boxes with score $\geq 0.7$. The count statistics
(Table~\ref{tab:sam3-thr}) show that a lower threshold trades missing objects for extra boxes, but not which side
of the trade the language model prefers: every box is given with its score and the model is told to trust a box
according to it, so it may tolerate low-score boxes. We evaluate the same list (prompts = the names in the
scene's questions, ``room door'' for door) at thresholds 0.1, 0.3, 0.5, 0.7 and 0.9 on all 5,130 questions with
the same model.
\emph{Result:} the threshold hardly matters (Table~\ref{tab:sam3-q6}). Over a threefold change in the
number of boxes, the full score stays between 68.3 and 69.2 and the debiased score between 62.2 and 62.8; 0.7 is
best or within 0.2 of the best on both. The one clear trend is counting, which falls steadily as the threshold
is lowered (74.0 at 0.9, 69.7 at 0.1): extra boxes are counted. Relative distance (66.3 to 68.9), absolute distance and
route planning change without a trend. No threshold reaches the video-only score (69.4 / 63.7), so the
threshold is not the lever; lowering it is useful only if a later stage removes the extra boxes.

\begin{table}[t]
\centering
\caption{VSI-Bench score for five box-score thresholds (question names + ``room door''). Boxes: mean number of
boxes in the list of a scene before the list is restricted to the objects a question names.}
\label{tab:sam3-q6}
\begin{tabular}{rrrrrrrr}
\toprule
Threshold & Boxes & Full & Debiased & Counting & Abs.\ distance & Rel.\ distance & Route planning \\
\midrule
0.1 & 32.6 & 68.3 & 62.2 & 69.7 & 56.8 & 68.0 & 52.1 \\
0.3 & 26.4 & 68.7 & 62.7 & 70.9 & 56.8 & 67.7 & 53.1 \\
0.5 & 19.5 & 68.5 & 62.2 & 72.2 & 57.3 & 66.3 & 50.5 \\
0.7 & 15.9 & 69.0 & 62.8 & 73.5 & 57.4 & 66.9 & 52.1 \\
0.9 & 10.0 & 69.2 & 62.2 & 74.0 & 56.7 & 68.9 & 52.1 \\
\bottomrule
\end{tabular}
\end{table}

\paragraph{Q7. The estimated scale is too large on average. Can it be improved, and should VGG-T-Det refine it?}
VGGT-SLAM reconstructs a scene up to an unknown scale. The pipeline estimates the metric scale from the objects:
for every well-seen instance of a class with a size prior, the ratio of the prior footprint and height to the
measured ones is taken, and the scale is the 35th percentile of these ratios (a low percentile, because partly
seen objects measure too small and would otherwise inflate the scale).
\emph{Known so far:} against the true scale (from the ground-truth camera poses), the estimate is too large in
the median by 4\% and in the mean by 11\% (Table~\ref{tab:sam3-q7}). The error is skewed: a tenth of the scenes
is more than a third too large, a tenth more than 13\% too small. Only 48\% of the scenes are within 10\% and
74\% within 20\%. The scale error is the largest geometric error of the list: in the ablation of
Section~\ref{sec:sam3}, ground-truth box sizes are the only geometric correction that helps the language model,
and in the query check for the detector the estimated scale costs 9 to 14 points of box coverage compared with
the true scale.

\begin{table}[t]
\centering
\caption{Estimated scale divided by true scale, per scene (question names + ``room door'').}
\label{tab:sam3-q7}
\begin{tabular}{lrrrrrr}
\toprule
Dataset & Scenes & Median & Mean & 10th--90th percentile & Within 10\% & Within 20\% \\
\midrule
ScanNet     &  88 & 1.06 & 1.14 & 0.89--1.33 & 49\% & 78\% \\
ARKitScenes & 150 & 1.05 & 1.10 & 0.86--1.35 & 43\% & 70\% \\
ScanNet++   &  49 & 1.03 & 1.07 & 0.83--1.30 & 65\% & 78\% \\
All         & 287 & 1.04 & 1.11 & 0.87--1.34 & 48\% & 74\% \\
\bottomrule
\end{tabular}
\end{table}

Ways to improve the estimate (the first three are tested below):
\begin{itemize}
\item \emph{Camera height.} The height of the camera above the floor is measured in the reconstruction and is
      close to a constant for hand-held video; it gives a second, object-free scale estimate. It can be combined
      with the object ratios or used to reject outliers.
\item \emph{More reliable objects.} Use only instances seen completely (mask not cut by the image border, many
      views), prefer classes with a narrow size distribution (door height, toilet, bed length) over classes with
      a wide one (table, cabinet), and weight the ratios by the spread of the class prior.
\item \emph{Calibrate the percentile.} The 35th percentile was set by hand; the bias in Table~\ref{tab:sam3-q7}
      suggests a lower one, and it can be fitted on the training scenes, per dataset if needed.
\item \emph{A learned metric prior.} A monocular metric-depth model on a few keyframes gives the scale without
      object priors.
\end{itemize}

\emph{Tested without training} (Table~\ref{tab:sam3-q7b}; constants fitted on 939 ScanNet training scenes, evaluated on
the 287 VSI-Bench scenes; the ground truth is used for the evaluation only).
The bias is easy to remove: taking the 30th instead of the 35th percentile of the object ratios centres the
estimate (median 1.01) and halves the share of scenes that are more than 30\% too large (14\% to 9\%). The spread
hardly changes: the typical error stays near 15\% and only half of the scenes are within 10\%. The other ideas
do not help. The camera height is not a usable constant: its true value is 0.8\,m in ScanNet and 1.3\,m in
ARKitScenes and ScanNet++, and it varies from 0.5 to 1.4\,m within ScanNet, so a scale from the camera height
alone is far worse, and mixing it with the object estimate is worse than the object estimate alone. Weighting
the classes by how reliable their size is, with the per-class bias removed, gives the same error as the plain
percentile. The limit is therefore the size priors themselves: object sizes vary too much within a class to fix
the scale to better than about 15\% from a handful of objects. A better scale needs another source of metric
information, a metric-depth model or the detector.

\begin{table}[t]
\centering
\caption{Scale estimators that use no ground truth. Error: mean absolute log ratio of estimated to true scale.}
\label{tab:sam3-q7b}
\begin{tabular}{lrrrrr}
\toprule
Estimator & Median ratio & Error & Within 10\% & Within 20\% & More than 30\% too large \\
\midrule
Object ratios, 35th percentile (current)          & 1.04 & 0.154 & 48\% & 74\% & 14\% \\
Object ratios, 30th percentile                    & 1.01 & 0.146 & 51\% & 76\% &  9\% \\
Object ratios, 25th percentile                    & 0.98 & 0.152 & 49\% & 76\% &  7\% \\
Object heights only, 35th percentile              & 1.06 & 0.202 & 48\% & 71\% & 18\% \\
Classes weighted by reliability, bias removed     & 0.95 & 0.155 & 45\% & 75\% &  8\% \\
Camera height only                                & 0.90 & 0.367 & 14\% & 40\% & 31\% \\
Object estimate and camera height, equal weights  & 0.96 & 0.216 & 28\% & 61\% & 20\% \\
\bottomrule
\end{tabular}
\end{table}

Should VGG-T-Det refine the scale? Today the detector cannot: it is given a ground-truth scale (the mean distance
of the scan points from the first camera) with its input, in training and at test time, so it never sees the
scale error. Two designs are possible. (i) The detector keeps receiving a scale, now the estimated one, and its
size regression has to absorb the error; this needs training with the estimated scale, or with the true scale
perturbed by the error distribution of Table~\ref{tab:sam3-q7}. (ii) The detector predicts one scale correction
per scene from all its queries together, which is applied to every box; a per-scene quantity is better
estimated from all objects at once than box by box. In both cases the question is whether the images and VGGT
tokens carry enough metric information beyond the object size priors the pipeline already uses.


\subsection{Answers so far to Q1, Q2 and Q5}
\label{sec:sam3-answers}

\paragraph{Q1: coverage of the detector's training data.}
Table~\ref{tab:sam3-a1} counts the annotated objects that fall into the 65 classes. A third (ScanNet) to more
than half (ScanNet++) of the annotated objects are outside them; these are mostly small clutter (book, paper,
box, bottle, socket, clothes, bag, light switch), spread over 3,126 distinct labels in ScanNet++. Only
ScanNet++ has boxes for all 65 classes. Seven classes exist in ScanNet++ alone: ceiling light (2,501 boxes),
power strip (241), mattress (138), cutting board (113), paper bag (110), coat rack (71) and shoe rack (49).
ARKitScenes has no boxes for 14 classes, among them counter, desk, monitor, closet and computer tower.
Whether SAM 3 finds the annotated objects in the training scenes is not measured yet.

\begin{table}[t]
\centering
\caption{Annotated objects of the three training datasets and the share that falls into the 65 classes.
Structure: walls, floors, ceilings and the generic ``object''.}
\label{tab:sam3-a1}
\begin{tabular}{lrrrrr}
\toprule
Dataset & Annotations & Structure & Objects & In the 65 classes & Classes with boxes \\
\midrule
ScanNet (EmbodiedScan, 1,201 scenes)     &  48,964 & 30.1\% & 34,220 & 66.4\% & 57 \\
ARKitScenes (EmbodiedScan, 4,498 scenes) & 102,623 & 16.3\% & 85,926 & 86.9\% & 51 \\
ScanNet++ (905 scenes)                   & 100,799 & 29.9\% & 70,672 & 42.4\% & 65 \\
\bottomrule
\end{tabular}
\end{table}

\paragraph{Q2: coverage of the Reason-190K questions.}
The objects a training question is about are the labels of its object list (Table~\ref{tab:sam3-a2};
room-size questions name none and are left out). Questions from ScanNet and ARKitScenes are covered almost
completely; the only notable gap is bicycle (287 ScanNet questions). Of the ScanNet++ questions, 28\% name an
object outside the 65 classes, most often light switch (2,337 questions), board (2,089), bottle, box, switch,
bag, slippers, book, blinds and jacket. That is about 18,400 of the 143,437 training questions (13\%).

\begin{table}[t]
\centering
\caption{Reason-190K questions whose named objects are all in the 65 classes. Exact: the class names and their
VSI-Bench synonyms. With renamings: also the dataset labels mapped to a class (ceiling lamp, shelf, mouse,
coat hanger, shoe, poster, office chair, couch).}
\label{tab:sam3-a2}
\begin{tabular}{lrrrr}
\toprule
Source & Questions & Exact & With renamings & Object mentions covered \\
\midrule
ScanNet     &  49,569 & 98.7\% & 98.9\% & 99.5\% \\
ARKitScenes &  28,485 & 99.1\% & 99.6\% & 99.8\% \\
ScanNet++   &  65,383 & 64.4\% & 71.8\% & 87.0\% \\
All         & 143,437 & 83.1\% & 86.7\% & 93.6\% \\
\bottomrule
\end{tabular}
\end{table}

\paragraph{Q5, wording: the proposed rewordings mostly do not help.}
Each rewording of Table~\ref{tab:sam3-q5-few} was run as a single SAM 3 prompt on the scenes whose questions
name the object and substituted for the benchmark name (Table~\ref{tab:sam3-a5}). Only one is a clear gain:
prompting ``radiator'' also where the question says ``heater'' raises the boxes from 6 to 21 for 46 objects and
the objects found from 11\% to 39\%. ``Ceiling lamp'', ``desktop computer'', ``PC case'' and ``extension cord with sockets'' find
fewer objects than the benchmark name. ``Kitchen counter'' and ``bag'' find nearly all objects (85\% and 88\%)
but add many extra boxes (14 for 13 counters, 22 for 43 backpacks), so they exchange one count error for the
other. ``Ventilation fan'' changes nothing. The benchmark names are thus mostly as good as or better than the alternatives, and the low recall of
ceiling light, computer tower, fan and power strip is not a wording problem.

\begin{table}[t]
\centering
\caption{Rewordings tested (box score $\geq 0.7$). Found: objects with a box of that name within 0.75\,m.
Exact: scenes with the right number of boxes.}
\label{tab:sam3-a5}
\begin{tabular}{llrrrrrr}
\toprule
Object (ground truth) & Prompt & Boxes & Extra & Missing & Found & Exact & Verdict \\
\midrule
heater / radiator (46) & benchmark name            &  6 &  0 & 40 & 11\% & 18\% & \\
                       & ``radiator''              & 21 &  1 & 26 & 39\% & 44\% & better \\
                       & ``wall radiator''         & 15 &  1 & 32 & 28\% & 35\% & better \\
ceiling light (116)    & benchmark name            & 53 &  1 & 64 & 39\% & 24\% & \\
                       & ``ceiling lamp''          & 21 &  0 & 95 & 16\% & 10\% & worse \\
computer tower (42)    & benchmark name            & 25 &  1 & 18 & 50\% & 14\% & \\
                       & ``desktop computer''      & 37 & 10 & 15 & 38\% & 14\% & worse \\
                       & ``PC case''               & 13 &  0 & 29 & 26\% &  0\% & worse \\
fan / exhaust fan (14) & benchmark name            &  5 &  1 & 10 & 21\% & 30\% & \\
                       & ``ventilation fan''       &  6 &  1 &  9 & 29\% & 30\% & no change \\
power strip (10)       & benchmark name            &  6 &  1 &  5 & 40\% & 29\% & \\
                       & ``extension cord with sockets'' & 0 & 0 & 10 & 0\% & 0\% & worse \\
counter (13)           & benchmark name            & 12 &  6 &  7 & 31\% & 18\% & \\
                       & ``kitchen counter''       & 27 & 14 &  0 & 85\% & 36\% & more found, more extra \\
backpack (43)          & benchmark name            & 29 &  0 & 14 & 67\% & 52\% & \\
                       & ``bag''                   & 64 & 22 &  1 & 88\% & 55\% & more found, more extra \\
                       & ``backpack'' and ``bag''  & 68 & 25 &  0 & 93\% & 52\% & more found, more extra \\
\bottomrule
\end{tabular}
\end{table}

\paragraph{Q5, decision.}
``Radiator'' is adopted: every scene whose questions say ``heater'' or ``radiator'' is prompted with ``radiator'',
and the box keeps the name the question uses. Over all 287 scenes this lowers the missing boxes of the list from
641 to 627 and leaves the extra boxes unchanged (414 to 415); the share of (scene, name) pairs with the exact
count rises from 54.2\% to 54.8\%. Radiators remain a weak class: 26 of 46 are still missing.
``Kitchen counter'' and ``bag'' are on hold: they find nearly all objects but add many extra boxes, so they are
useful only once a later stage removes extra boxes. The other rewordings are rejected. The prompt set is
therefore the names of the scene's questions with two substitutions, ``room door'' for door and ``radiator''
for heater.

\paragraph{Q5, extra boxes: half of them are on nothing that is annotated.}
Of 2,967 boxes with score $\geq 0.7$ in all scenes, 2,368 match a ground-truth object of their name and 599 do
not. Of these 599, 299 (50\%) lie where nothing is annotated, 163 (27\%) duplicate an object that already has a
box, and 137 (23\%) lie on an object of another class (Table~\ref{tab:sam3-a5x}). This corrects two guesses of
Table~\ref{tab:sam3-q5-many}. Duplicates are not the main cause for table, chair and sofa, so better merging
would remove about a quarter of the extra boxes at most; only for bed do duplicates lead. And tv boxes are
rarely on monitors (2 cases). Chair and sofa are confused in both directions, and pillow and cushion boxes often
land on the sofa or bed they lie on. Whether a box on nothing annotated is a real object without annotation or a
false detection cannot be told from the counts.

\begin{table}[t]
\centering
\caption{What the boxes of a name are (all 287 scenes, score $\geq 0.7$).}
\label{tab:sam3-a5x}
\begin{tabular}{lrrrrrl}
\toprule
Name & Boxes & Matched & Duplicate & Other class & Nothing annotated & Main other class \\
\midrule
table   & 454 & 317 & 48 & 25 & 64 & chair 11, bed 5 \\
chair   & 864 & 739 & 23 & 27 & 75 & sofa 17, table 6 \\
sofa    & 144 & 106 & 11 & 10 & 17 & chair 8 \\
door    & 163 & 131 &  6 &  2 & 24 & -- \\
bed     & 106 &  77 & 14 &  2 & 13 & -- \\
pillow  &  93 &  70 &  9 & 13 &  1 & sofa 6, bed 4 \\
tv      &  96 &  78 &  8 &  5 &  5 & monitor 2 \\
cushion &  12 &   5 &  1 &  6 &  0 & sofa 3, chair 3 \\
\midrule
All names & 2,967 & 2,368 & 163 & 137 & 299 & \\
\bottomrule
\end{tabular}
\end{table}

\subsection{True scale against estimated scale}
\label{sec:sam3-scale}

The SAM 3 + VGGT-SLAM list estimates its metric scale from object sizes (Q7). To measure what that estimate costs, the
same list (prompts from the scene's questions, ``room door'', ``radiator'') was scored once as it is and once with
its scale replaced by the true one, obtained by fitting the reconstructed cameras to the ground-truth cameras.
Nothing else changes: the same boxes, only rescaled about the first camera. The estimated scale is a median 10\%
off (49\% of the scenes within 10\%, 87\% within 30\%).

\paragraph{3D boxes.} Table~\ref{tab:sam3-scale-3d}: predicted boxes are matched to ground-truth boxes of the same
label (287 VSI-Bench scenes, 3,194 objects); precision counts boxes with score $\geq 0.7$. The true scale improves
every loose-match measure by 5 to 9 points, on all three datasets (Table~\ref{tab:sam3-scale-ds}). Tight matches do
not improve (mAP@0.50 stays near 4): the scale only fixes the distance of every box from the camera, not the boxes'
shape and placement errors.

\begin{table}[t]
\centering
\caption{SAM 3 + VGGT-SLAM boxes with the estimated and the true scale (287 VSI-Bench scenes, 3,194 objects).}
\label{tab:sam3-scale-3d}
\begin{tabular}{lrrr}
\toprule
 & Estimated scale & True scale & Change \\
\midrule
Recall, IoU $\geq 0.25$             & 30.5 & 38.2 & $+7.7$ \\
Precision, IoU $\geq 0.25$          & 32.5 & 40.3 & $+7.8$ \\
mAP@0.25                            & 17.2 & 22.2 & $+5.0$ \\
Recall, IoU $\geq 0.5$              &  8.1 & 11.2 & $+3.1$ \\
mAP@0.50                            &  4.4 &  4.0 & $-0.4$ \\
Recall, centre within 0.5\,m        & 67.3 & 75.9 & $+8.6$ \\
Precision, centre within 0.5\,m     & 65.8 & 73.9 & $+8.1$ \\
\bottomrule
\end{tabular}
\end{table}

\begin{table}[t]
\centering
\caption{Recall at IoU $\geq 0.25$ by dataset.}
\label{tab:sam3-scale-ds}
\begin{tabular}{lrrr}
\toprule
Dataset & Objects & Estimated scale & True scale \\
\midrule
ScanNet     & 1,224 & 32.0 & 41.9 \\
ARKitScenes &   854 & 34.8 & 43.0 \\
ScanNet++   & 1,116 & 25.4 & 30.4 \\
\bottomrule
\end{tabular}
\end{table}

\paragraph{VSI-Bench.} Table~\ref{tab:sam3-scale-vsi} gives the score of the default language model (mix-v2-1) with
the two lists. The true scale helps the two metric tasks a little (absolute distance $+1.0$, size $+1.3$) and leaves
the overall score unchanged; counting even drops by 1.1, within the noise of a single run. The scale error therefore
costs far less than the errors in the list itself: which objects are listed and how loose their boxes are.

\begin{table}[t]
\centering
\caption{VSI-Bench (all 5,130 questions) with the SAM 3 + VGGT-SLAM list at the estimated and the true scale.}
\label{tab:sam3-scale-vsi}
\begin{tabular}{lrrrrr}
\toprule
SAM 3 list & Full & Debiased & Abs.\ distance & Size & Counting \\
\midrule
Estimated scale & 69.0 & 62.9 & 57.5 & 73.4 & 73.5 \\
True scale      & 69.1 & 63.0 & 58.5 & 74.7 & 72.4 \\
\bottomrule
\end{tabular}
\end{table}

\paragraph{Which results use the true scale.} Every VGG-T-Det variant (the released 18-class model, cls-31, cls-63,
cls-65 and IndoorDet-65-31k) multiplies VGGT's normalized point cloud by the mean distance of the ground-truth scan
from the first camera and uses the first camera's ground-truth pose, in training and at test time; the proposals of
IndoorDet-65-31k are placed with the true scale as well. All object lists built from a detector, and all language-model
training data (ground-truth and detector boxes), are therefore metric with the true scale. Only the SAM 3 + VGGT-SLAM
lists use an estimated scale, so they are the only video-only lists, and the language model sees a scale error only
at test time with them.

Table~\ref{tab:sam3-scale-use} lists which runs use the true scale.

\begin{table}[t]
\centering
\caption{Runs with and without the true scale. ``True scale'' = metric scale from the ground-truth scan (detectors) or
from fitting the reconstructed cameras to the ground-truth cameras (SAM 3 lists); ``estimated'' = from object size
priors, with up direction and walls from the reconstruction and the first camera as origin (no ground truth).}
\label{tab:sam3-scale-use}
\small
\begin{tabular}{llll}
\toprule
Run & Stage & Scale & Used for \\
\midrule
\multicolumn{4}{l}{\emph{Without the true scale (video only)}} \\
SAM 3 + VGGT-SLAM lists, all prompt variants (fixed 65 names, question & test & estimated & VSI-Bench, e.g.\ 69.0 / 62.9 \\
\quad names, ``room door'', ``radiator'', route phrases, thresholds 0.1--0.9) & & & \\
VGGT-SLAM lists with oracle 2D boxes (\texttt{slam\_manh}, \texttt{slamall\_manh}) & test & estimated & VSI-Bench \\
SAM 3 list ablation: predicted centre and size cells & test & estimated & VSI-Bench (Section~\ref{sec:sam3}) \\
SAM 3 list, estimated-scale arm of Table~\ref{tab:sam3-scale-3d} & test & estimated & 3D evaluation, VSI-Bench \\
Video-only model & train, test & no boxes & VSI-Bench 69.4 / 63.7 \\
\midrule
\multicolumn{4}{l}{\emph{With the true scale}} \\
VGG-T-Det: released 18-class, ES-56, cls-31, VSI-40, cls-63, cls-65 & train, test & true & detection; VSI-Bench lists \\
IndoorDet-65-31k (test2-hybrid), and its SAM 3 proposals & train, test & true & detection; VSI-Bench lists \\
Lists from a detector (cls-31; IndoorDet-65-31k refine and & test & true & VSI-Bench, e.g.\ 69.9 / 62.7, 69.8 / 63.2 \\
\quad threshold lists) & & & \\
VGGT-SLAM lists \texttt{\_truescale}; SAM 3 list, true-scale arm & test & true & VSI-Bench, 3D evaluation \\
Language-model SFT data (mix-v2-0, mix-v2-1, Exp1): ground-truth & train & true & every SFT run \\
\quad and detector boxes & & & \\
\bottomrule
\end{tabular}
\end{table}
; needs \usepackage{booktabs}.
% All numbers: VSI-Bench, 5,130 questions, 288 scenes (287 processed); box score threshold 0.7 unless stated.
% Sources: VGGT-SLAM/exps/vsibench_objects (s3_alt_stats.py, s3_badlist_thr.py, qvocab.py, s3_rows.py),
% 3d-scene-lang/runs/eval_q35vs_gandalf/{mixv2_1_sam3_full,mixv2_1_s3q_roomdoor}_rvs_question.

\section{Open-vocabulary object lists with SAM 3}
\label{sec:sam3}

\paragraph{Pipeline.}
SAM 3 is run on up to 100 VGGT-SLAM keyframes per video with one text prompt per object name
(per-frame confidence $\geq 0.5$). Each mask is lifted to 3D with the SLAM points, masks of the same name are
merged across frames, and each merged instance gives one box. The box score is the mean of the three best
SAM 3 scores times $1-\exp(-\mathrm{views}/2)$. No ground truth is used. The resulting list of
\{label, box, score\} entries is given to the language model together with the video.

\paragraph{The object list, not the geometry, limits the benefit.}
Replacing the predicted centres and sizes by ground truth while keeping the SAM 3 list raises the VSI-Bench
score only from 68.6 to 69.9, whereas the full ground-truth list reaches 77.3. The errors that matter are
extra boxes and missing objects. With question-name prompts and ``room door'' (below), the list has 414 extra
and 641 missing boxes against 3,194 ground-truth objects of the names the questions use.

\subsection{Run time}
\label{sec:sam3-time}

Table~\ref{tab:sam3-time} gives the cost of the two stages, measured on 597 ScanNet++ scenes processed together
on four NVIDIA RTX PRO 6000 Blackwell Max-Q Workstation Edition GPUs (96\,GB each, machine gandalf) with three
workers per GPU; the GPUs ran nothing else. VGGT-SLAM receives up to 600 frames per scene and
keeps 100 keyframes; SAM 3 runs all 65 class prompts on each of the 100 keyframes (per-frame confidence 0.3).
SAM 3 costs about eight times as much as VGGT-SLAM per scene. Its time hardly varies between scenes, since every
prompt is run on every keyframe (about 4\,s per keyframe for 65 prompts); with three workers sharing a GPU it is
limited by GPU compute, so more workers per GPU do not raise the throughput. For comparison, SAM 3 on ARKitScenes on
four NVIDIA RTX 6000 Ada GPUs (46\,GB each, machine wukong; two workers per GPU) took about
250\,s per scene per worker with 51 prompts.

\begin{table}[t]
\centering
\caption{Run time per scene of the two stages (597 ScanNet++ scenes, 4 NVIDIA RTX PRO 6000 Blackwell Max-Q GPUs,
96\,GB, 12 workers).}
\label{tab:sam3-time}
\begin{tabular}{lrrrr}
\toprule
Stage & \multicolumn{3}{c}{Time per scene per worker} & Throughput \\
      & Median & Mean & 90th percentile & (12 workers) \\
\midrule
VGGT-SLAM (600 frames in, 100 keyframes)   &  48\,s &  48\,s &  74\,s & 10 scenes per minute \\
SAM 3 (65 prompts on 100 keyframes)        & 394\,s & 392\,s & 398\,s & 110 scenes per hour \\
\bottomrule
\end{tabular}
\end{table}

\subsection{Prompt wording changes the number of boxes}

\paragraph{``door'' versus ``room door''.}
The prompt ``door'' also fires on cabinet, wardrobe and shower doors. It alone produced 37\% of all extra boxes
(231 of 620). Prompting ``room door'' instead removes most of them at the cost of some recall
(Table~\ref{tab:sam3-door}).

\begin{table}[t]
\centering
\caption{Door boxes for the two prompts (170 ground-truth doors in 114 scenes).
Found: a box of that label within 0.75\,m of the ground-truth centre.}
\label{tab:sam3-door}
\begin{tabular}{lrrrrr}
\toprule
Prompt & Boxes & Extra & Missing & Scenes with exact count & Found \\
\midrule
``door''      & 391 & 231 & 10 & 31\% & 85\% \\
``room door'' & 163 &  25 & 32 & 64\% & 76\% \\
\bottomrule
\end{tabular}
\end{table}

\paragraph{A more specific word is not always better.}
For tables, every rewording we tried is worse than the plain name (Table~\ref{tab:sam3-table}): ``desk'' removes
13 extra boxes but loses 61 more tables. The extra table boxes are duplicates of real tables and tables without
annotation, so they call for better merging, not for another word.

\begin{table}[t]
\centering
\caption{Table boxes for four prompts (406 ground-truth tables, 285 scenes, all with ``room door'').}
\label{tab:sam3-table}
\begin{tabular}{lrrrrr}
\toprule
Prompt & Boxes & Extra & Missing & Scenes with exact count & Found \\
\midrule
``table'' (benchmark name)      & 450 & 101 &  57 & 41\% & 75\% \\
``desk''                        & 376 &  88 & 118 & 35\% & 62\% \\
``table or desk''               & 529 & 162 &  39 & 34\% & 79\% \\
``dining table or office desk'' & 287 &  58 & 177 & 36\% & 50\% \\
\bottomrule
\end{tabular}
\end{table}

\paragraph{Which names are detected too rarely or too often.}
Table~\ref{tab:sam3-names} lists the names with the largest count errors. Two groups appear.
Names detected too rarely are mostly ones whose benchmark wording is unusual for the object
(``heater'' for a wall radiator, ``ceiling light'', ``computer tower'') or small objects (basket, shoes, bucket,
pan, kettle), which have no extra boxes at all. Names detected too often (table, chair, sofa, bed) are large
furniture that is split into several boxes or duplicated. ``pillow'' and ``cushion'' fire on the same objects.
Only the door and table wordings have been tested; the others are candidates.

\begin{table}[t]
\centering
\caption{Count errors per object name (question-name prompts with ``room door''). Scenes: those whose questions
use the name and that contain the object.}
\label{tab:sam3-names}
\begin{tabular}{lrrrrr}
\toprule
Name & Ground truth & Boxes & Extra & Missing & Scenes with too few \\
\midrule
\multicolumn{6}{l}{\emph{Too few boxes}} \\
heater / radiator & 46 & 6 & 0 & 40 & 82\% \\
ceiling light     & 116 & 53 & 1 & 64 & 72\% \\
computer tower    & 42 & 25 & 1 & 18 & 71\% \\
fan / exhaust fan & 14 & 5 & 1 & 10 & 60\% \\
power strip       & 10 & 6 & 1 & 5 & 57\% \\
backpack          & 43 & 29 & 0 & 14 & 48\% \\
window            & 83 & 68 & 11 & 26 & 40\% \\
\midrule
\multicolumn{6}{l}{\emph{Too many boxes}} \\
table & 409 & 454 & 102 & 57 & 23\% \\
chair & 897 & 864 & 73 & 106 & 26\% \\
sofa  & 146 & 144 & 30 & 32 & 31\% \\
bed   & 82 & 106 & 26 & 2 & 3\% \\
door (``room door'') & 170 & 163 & 25 & 32 & 18\% \\
\bottomrule
\end{tabular}
\end{table}

\paragraph{The score threshold trades missing objects for extra ones.}
Lowering the box-score threshold does not fix the list (Table~\ref{tab:sam3-thr}, fixed 65-name prompts): at 0.3
fewer labels are short of boxes, but there is about one extra box per ground-truth object. About 7\% of labels
are never found at any threshold, because SAM 3 gives no mask above its per-frame confidence.

\begin{table}[t]
\centering
\caption{Per (scene, label) count against ground truth for four score thresholds.}
\label{tab:sam3-thr}
\begin{tabular}{rrrrr}
\toprule
Threshold & Exact & Too few & Too many & Extra boxes per ground-truth object \\
\midrule
0.7 & 44.4\% & 22.6\% & 33.0\% & 0.43 \\
0.5 & 41.6\% & 15.6\% & 42.8\% & 0.64 \\
0.3 & 30.9\% &  9.9\% & 59.2\% & 1.06 \\
0.2 & 24.2\% &  7.0\% & 68.8\% & 1.48 \\
\bottomrule
\end{tabular}
\end{table}

\subsection{The vocabulary of the questions}

\paragraph{Seven of eight tasks use a closed list of 65 names.}
Outside route planning, every object in a VSI-Bench question is one of 65 names, and the wording depends on the
source dataset (``trash can'' / ``trash bin'', ``heater'' / ``radiator'', ``whiteboard'' / ``blackboard'').
Prompting SAM 3 with the names a scene's own questions use, in their original wording, gives 15.9 boxes per
question instead of 41.9 for a fixed list of all names.

\paragraph{Route-planning questions are open-vocabulary.}
The 194 route-planning questions are free text. Of their 979 object mentions (start, facing direction, goal and
every ``go forward until''), only 505 (52\%) are one of the 65 names. The rest are other objects
(sink in 28 scenes, cabinet 11, bathroom sink 8, desk 6, doorframe 6, wardrobe 4),
objects with an attribute (``blue chair'', ``open door'', ``glass coffee table''),
spelling variants (``trashbin'', ``bathtube'', ``refridgerator'', ``dinning table'')
and places that are not objects (wall, ``into the room'', ``outside'').
A fixed vocabulary therefore gives no box at all for about half of the objects these questions refer to.

\paragraph{Extraction rule.}
We take each mention in the question's own words, remove articles, position words and trailing relations
(``the sink on the left'' $\to$ ``sink'', ``foot of the bed'' $\to$ ``bed''), keep attributes
(``blue chair''), and drop walls, rooms and ``outside''. This gives 183 distinct phrases; 241 (scene, phrase)
pairs in 147 scenes are not already question names and become additional prompts. Only a plain ``door'' is
prompted as ``room door''; ``doorframe'', ``doorway'', ``open door'' or ``kitchen door'' are used as written.
When a route phrase and a benchmark name find the same object, the list keeps one box, shown under the name the
question uses.

\subsection{Effect on VSI-Bench}

Table~\ref{tab:sam3-vsi} compares the lists with the same language model (Qwen3.5-4B, mix-v2-1).
Question-name prompts with ``room door'' gain 0.4 / 0.5 points over the fixed 65-name prompts. The gain is in
counting ($+1.5$) and absolute distance ($+0.8$), the tasks most affected by extra door boxes; the other tasks
change by at most 0.6. The list still does not beat the video alone.
Adding the route phrases changes the boxes shown for 130 of the 194 route-planning questions, yet the
route-planning score stays at 52.1: only 10 answers change, 4 from wrong to right and 4 from right to wrong.
Section~\ref{sec:sam3-route} examines why.

\begin{table}[t]
\centering
\caption{VSI-Bench score (mean over tasks) on all 5,130 questions and on the debiased subset.}
\label{tab:sam3-vsi}
\begin{tabular}{lrrrrr}
\toprule
Object list & Full & Debiased & Counting & Abs.\ distance & Route planning \\
\midrule
None (video only)                               & 69.4 & 63.7 & -- & -- & -- \\
SAM 3, fixed 65 names                           & 68.6 & 62.3 & 72.0 & 56.6 & 51.5 \\
SAM 3, question names + ``room door''           & 69.0 & 62.8 & 73.5 & 57.4 & 52.1 \\
SAM 3, question names + ``room door'' + route phrases & 69.0 & 62.7 & 73.3 & 57.4 & 52.1 \\
SAM 3 list with ground-truth centre and size    & 69.9 & 63.9 & -- & -- & -- \\
Ground-truth list                               & 77.3 & 73.9 & -- & -- & -- \\
\bottomrule
\end{tabular}
\end{table}

\subsection{Route planning: are the named objects in the list?}
\label{sec:sam3-route}

Route planning is the task an open vocabulary should help most, since half of its object mentions are outside
the 65 names. We therefore check, for each of the 194 questions, whether every object it names has a box in
the list that includes the route phrases. A question names 659 distinct objects in total, 418 by a benchmark
name and 241 by another phrase.

\begin{table}[t]
\centering
\caption{Route-planning questions: objects named in the question that have at least one box of that name,
for four box-score thresholds. The evaluated list uses 0.7.}
\label{tab:sam3-route}
\begin{tabular}{rrrr}
\toprule
Threshold & Questions with every object found & Benchmark names found & Other phrases found \\
\midrule
0.7  & 120 / 194 (62\%) & 375 / 418 (90\%) & 198 / 241 (82\%) \\
0.5  & 145 / 194 (75\%) & 394 / 418 (94\%) & 210 / 241 (87\%) \\
0.3  & 157 / 194 (81\%) & 400 / 418 (96\%) & 218 / 241 (90\%) \\
none & 164 / 194 (85\%) & 403 / 418 (96\%) & 222 / 241 (92\%) \\
\bottomrule
\end{tabular}
\end{table}

\paragraph{Coverage is incomplete.}
At the evaluated threshold only 62\% of the questions have a box for every object they name
(Table~\ref{tab:sam3-route}). The objects missing most often are heater (9 questions), door (7), window (4),
cabinet (3) and shower (3). Heater is also the main object never found at any threshold (8 questions),
in line with Table~\ref{tab:sam3-names}; the remaining never-found objects are single cases such as
``column'', ``trophies'' or ``paper shredder box''.

\paragraph{Found objects are often ambiguous.}
A route question refers to one specific object (``the door'', ``the blue chair''). At threshold 0.7,
149 of the 375 found benchmark names and 91 of the 198 found phrases have several boxes of that name, so the
list does not say which box is meant. Lowering the threshold closes part of the coverage gap but widens this
one: with no threshold, 295 of the 403 found benchmark names have several boxes.

\paragraph{The language model barely uses the added boxes.}
The route phrases change the boxes shown for 130 of the 194 questions, and the 6 questions that matched no box
before now match some. Yet only 10 answers change, 4 from wrong to right and 4 from right to wrong; 101 answers
are correct with and without the phrases (52.1\%). A likely reason, which we have not
tested, is that the model was fine-tuned on lists that use the dataset label names, so boxes under free phrases
are an input it has not been trained to use.

\paragraph{Limits of this check.}
``Found'' means that a box with the name exists. Free phrases have no ground-truth boxes, so we cannot verify
that the box is on the right object. Mentions that are not objects (walls, rooms, ``outside'') are not counted.

\paragraph{Conclusion for route planning.}
Three gaps act together: 38\% of the questions miss at least one named object, about 40\% of the found objects
are ambiguous, and the model does not exploit the additional boxes. Extending the vocabulary alone therefore
gives no gain (52.1 with and without route phrases); the task needs boxes that identify the referred instance
and a model trained on such lists.

\paragraph{Summary.}
(i) Text prompts are not interchangeable: one word (``room'') removes 206 of 231 extra door boxes, while a
plausible synonym for table loses 61 objects.
(ii) Benchmark wording should be kept per dataset, and ambiguous names disambiguated, not replaced by synonyms.
(iii) A closed vocabulary covers only about half of the objects in free-text questions; taking the prompts from
the question raises coverage to 82--90\% of the named objects, but does not change the route-planning score
(Section~\ref{sec:sam3-route}).
(iv) Wording fixes a minority of the errors: after them, missing objects (641) outnumber extra boxes (414), and
neither a lower threshold nor rewording removes both.

\subsection{Open questions}
\label{sec:sam3-open}

\paragraph{Q1. Do the 65 classes (with ``room door'') cover the detector's training data?}
The detector (VGG-T-Det) is trained on ScanNet, ScanNet++ and ARKitScenes. Do the 65 class names, prompted as
in Section~\ref{sec:sam3} with ``room door'' for door, cover the annotated objects of these three datasets, and
which annotated objects fall outside them?
\emph{Known so far:} the 65 classes are the VSI-Bench names, so every class has a benchmark use, but the
converse is not measured. In the 242 ScanNet++ scenes that are ready, 7,881 of 27,878 annotations map to one of
the 65 classes (the rest include walls, floors and the generic ``object''); seven classes have boxes only in
ScanNet++ (ceiling light, cutting board, paper bag, shoe rack, coat rack, power strip, mattress) and two have
none there (fireplace, piano). The shares for ScanNet and ARKitScenes, and whether SAM 3 with these prompts
finds the annotated objects of each class in the training scenes, are still to be measured.

\paragraph{Q2. Do the 65 classes cover the objects named in the Reason-190K questions?}
The language model is fine-tuned on Reason-190K. If its questions name objects outside the 65 classes, a
65-class list gives no box for them in training, as it does for route planning in VSI-Bench.
\emph{Known so far:} not measured for the 65-class set. The measurement is the one of
Section~\ref{sec:sam3-route} applied to the training questions: the share of questions whose named objects are
all among the 65 names, per source dataset.

\paragraph{Q3. Should the whole pipeline be open-vocabulary?}
If the answers to Q1 and Q2 show large gaps, the alternative is to drop the fixed class set: the label of a box
comes from the SAM 3 prompt, the classification head of the detector is removed, and the detector only refines
the centre and size of each box and decides whether to keep it. For this to work the detector must receive
more candidate boxes than are finally kept, so the SAM 3 score threshold would be lowered and the removal of
extra boxes left to the keep decision.
\emph{Known so far:} lowering the box-score threshold from 0.7 to 0.3 reduces the labels with too few boxes
from 22.6\% to 9.9\% but raises the extra boxes per ground-truth object from 0.43 to 1.06
(Table~\ref{tab:sam3-thr}), so a keep decision would have to remove about one box per object. About 7\% of
labels are found at no box-score threshold; recovering them needs a lower SAM 3 per-frame confidence (0.5 now),
which has not been tried. Two risks are open: a correct box on an object without annotation is trained as
``drop'', and the detector cannot add an object that SAM 3 missed.

\paragraph{Q4. Should SAM 3 run on overlapping video segments that are merged afterwards, like submaps?}
VGGT-SLAM handles long videos by reconstructing submaps of consecutive frames that share overlapping frames and
then aligning them. The same scheme could be used for the objects: run SAM 3 on each video segment, with the
overlapping frames in common, build the object instances per segment, and merge the segment instances at the end,
using the overlapping frames to decide which instances are the same object.
\emph{Known so far:} the current pipeline does not use segments. SAM 3 is applied as an image model to each
keyframe independently (at most 100 keyframes spread evenly over the video, so long videos are subsampled), and
masks are associated only afterwards in 3D, by the overlap of their lifted points over the whole video. Extra
boxes from this global merge are the main count error for large furniture (Table~\ref{tab:sam3-names}), and
errors of the global reconstruction enter every association. Open points: whether SAM 3's video mode, which
tracks an object through consecutive frames, gives cleaner instances within a segment than per-frame masks;
whether merging per segment and then across segments produces fewer duplicates than one global merge; and how
many keyframes per segment are needed so that small objects are seen in several views.

\paragraph{Q5. Which count errors can be fixed by wording, and which need better merging?}
Two groups of names have wrong box counts (question-name prompts with ``room door'', box score $\geq 0.7$).
For the names with too few boxes, the benchmark wording is the likely cause, and
Table~\ref{tab:sam3-q5-few} lists a wording to try for each. For the names with too many boxes, the likely
cause is the merging of masks into instances, not the wording (Table~\ref{tab:sam3-q5-many}).
\emph{Known so far:} only the door and table wordings have been tested
(Tables~\ref{tab:sam3-door} and~\ref{tab:sam3-table}). The wordings and causes in the two tables below are
hypotheses read from the counts; none has been tested. Small objects with too few boxes and no extra boxes
(basket, shoes, bucket, pan, kettle, blanket, bowl, cup) are probably not a wording problem; a lower SAM 3
per-frame confidence is the more likely remedy.

\begin{table}[t]
\centering
\caption{Names with too few boxes, where the wording is the likely cause, and the wording to try (untested).}
\label{tab:sam3-q5-few}
\begin{tabular}{lrrrl}
\toprule
Name & Ground truth & Boxes & Scenes with too few & Wording to try \\
\midrule
heater / radiator & 46  & 6  & 82\% & ``radiator'' for both, or ``wall radiator'' \\
ceiling light     & 116 & 53 & 72\% & ``ceiling lamp'' (the ScanNet++ annotators' word) \\
computer tower    & 42  & 25 & 71\% & ``desktop computer'', ``PC case'' \\
fan / exhaust fan & 14  & 5  & 60\% & ``ventilation fan'' \\
power strip       & 10  & 6  & 57\% & ``extension cord with sockets'' \\
counter           & 13  & 12 & 55\% & ``kitchen counter'' \\
backpack          & 43  & 29 & 48\% & ``bag'' as a second prompt \\
window            & 83  & 68 & 40\% & unclear; may be visibility, not wording \\
\bottomrule
\end{tabular}
\end{table}

\begin{table}[t]
\centering
\caption{Names with too many boxes, where merging rather than wording is the likely cause (untested).}
\label{tab:sam3-q5-many}
\begin{tabular}{lrrrl}
\toprule
Name & Ground truth & Boxes & Extra & Likely cause \\
\midrule
table            & 409 & 454 & 102 & duplicates, tables without annotation \\
chair            & 897 & 864 & 73  & duplicates (also 106 missing) \\
sofa             & 146 & 144 & 30  & sectional sofas split into pieces \\
bed              & 82  & 106 & 26  & bed split into several boxes \\
tv               & 94  & 96  & 15  & monitors detected as tv \\
pillow / cushion & 100 & 105 & 20  & the two names detect the same objects \\
\bottomrule
\end{tabular}
\end{table}

\paragraph{Q6. Which box-score threshold is best for the language model?}
All VSI-Bench results above use boxes with score $\geq 0.7$. The count statistics
(Table~\ref{tab:sam3-thr}) show that a lower threshold trades missing objects for extra boxes, but not which side
of the trade the language model prefers: every box is given with its score and the model is told to trust a box
according to it, so it may tolerate low-score boxes. We evaluate the same list (prompts = the names in the
scene's questions, ``room door'' for door) at thresholds 0.1, 0.3, 0.5, 0.7 and 0.9 on all 5,130 questions with
the same model.
\emph{Result:} the threshold hardly matters (Table~\ref{tab:sam3-q6}). Over a threefold change in the
number of boxes, the full score stays between 68.3 and 69.2 and the debiased score between 62.2 and 62.8; 0.7 is
best or within 0.2 of the best on both. The one clear trend is counting, which falls steadily as the threshold
is lowered (74.0 at 0.9, 69.7 at 0.1): extra boxes are counted. Relative distance (66.3 to 68.9), absolute distance and
route planning change without a trend. No threshold reaches the video-only score (69.4 / 63.7), so the
threshold is not the lever; lowering it is useful only if a later stage removes the extra boxes.

\begin{table}[t]
\centering
\caption{VSI-Bench score for five box-score thresholds (question names + ``room door''). Boxes: mean number of
boxes in the list of a scene before the list is restricted to the objects a question names.}
\label{tab:sam3-q6}
\begin{tabular}{rrrrrrrr}
\toprule
Threshold & Boxes & Full & Debiased & Counting & Abs.\ distance & Rel.\ distance & Route planning \\
\midrule
0.1 & 32.6 & 68.3 & 62.2 & 69.7 & 56.8 & 68.0 & 52.1 \\
0.3 & 26.4 & 68.7 & 62.7 & 70.9 & 56.8 & 67.7 & 53.1 \\
0.5 & 19.5 & 68.5 & 62.2 & 72.2 & 57.3 & 66.3 & 50.5 \\
0.7 & 15.9 & 69.0 & 62.8 & 73.5 & 57.4 & 66.9 & 52.1 \\
0.9 & 10.0 & 69.2 & 62.2 & 74.0 & 56.7 & 68.9 & 52.1 \\
\bottomrule
\end{tabular}
\end{table}

\paragraph{Q7. The estimated scale is too large on average. Can it be improved, and should VGG-T-Det refine it?}
VGGT-SLAM reconstructs a scene up to an unknown scale. The pipeline estimates the metric scale from the objects:
for every well-seen instance of a class with a size prior, the ratio of the prior footprint and height to the
measured ones is taken, and the scale is the 35th percentile of these ratios (a low percentile, because partly
seen objects measure too small and would otherwise inflate the scale).
\emph{Known so far:} against the true scale (from the ground-truth camera poses), the estimate is too large in
the median by 4\% and in the mean by 11\% (Table~\ref{tab:sam3-q7}). The error is skewed: a tenth of the scenes
is more than a third too large, a tenth more than 13\% too small. Only 48\% of the scenes are within 10\% and
74\% within 20\%. The scale error is the largest geometric error of the list: in the ablation of
Section~\ref{sec:sam3}, ground-truth box sizes are the only geometric correction that helps the language model,
and in the query check for the detector the estimated scale costs 9 to 14 points of box coverage compared with
the true scale.

\begin{table}[t]
\centering
\caption{Estimated scale divided by true scale, per scene (question names + ``room door'').}
\label{tab:sam3-q7}
\begin{tabular}{lrrrrrr}
\toprule
Dataset & Scenes & Median & Mean & 10th--90th percentile & Within 10\% & Within 20\% \\
\midrule
ScanNet     &  88 & 1.06 & 1.14 & 0.89--1.33 & 49\% & 78\% \\
ARKitScenes & 150 & 1.05 & 1.10 & 0.86--1.35 & 43\% & 70\% \\
ScanNet++   &  49 & 1.03 & 1.07 & 0.83--1.30 & 65\% & 78\% \\
All         & 287 & 1.04 & 1.11 & 0.87--1.34 & 48\% & 74\% \\
\bottomrule
\end{tabular}
\end{table}

Ways to improve the estimate (the first three are tested below):
\begin{itemize}
\item \emph{Camera height.} The height of the camera above the floor is measured in the reconstruction and is
      close to a constant for hand-held video; it gives a second, object-free scale estimate. It can be combined
      with the object ratios or used to reject outliers.
\item \emph{More reliable objects.} Use only instances seen completely (mask not cut by the image border, many
      views), prefer classes with a narrow size distribution (door height, toilet, bed length) over classes with
      a wide one (table, cabinet), and weight the ratios by the spread of the class prior.
\item \emph{Calibrate the percentile.} The 35th percentile was set by hand; the bias in Table~\ref{tab:sam3-q7}
      suggests a lower one, and it can be fitted on the training scenes, per dataset if needed.
\item \emph{A learned metric prior.} A monocular metric-depth model on a few keyframes gives the scale without
      object priors.
\end{itemize}

\emph{Tested without training} (Table~\ref{tab:sam3-q7b}; constants fitted on 939 ScanNet training scenes, evaluated on
the 287 VSI-Bench scenes; the ground truth is used for the evaluation only).
The bias is easy to remove: taking the 30th instead of the 35th percentile of the object ratios centres the
estimate (median 1.01) and halves the share of scenes that are more than 30\% too large (14\% to 9\%). The spread
hardly changes: the typical error stays near 15\% and only half of the scenes are within 10\%. The other ideas
do not help. The camera height is not a usable constant: its true value is 0.8\,m in ScanNet and 1.3\,m in
ARKitScenes and ScanNet++, and it varies from 0.5 to 1.4\,m within ScanNet, so a scale from the camera height
alone is far worse, and mixing it with the object estimate is worse than the object estimate alone. Weighting
the classes by how reliable their size is, with the per-class bias removed, gives the same error as the plain
percentile. The limit is therefore the size priors themselves: object sizes vary too much within a class to fix
the scale to better than about 15\% from a handful of objects. A better scale needs another source of metric
information, a metric-depth model or the detector.

\begin{table}[t]
\centering
\caption{Scale estimators that use no ground truth. Error: mean absolute log ratio of estimated to true scale.}
\label{tab:sam3-q7b}
\begin{tabular}{lrrrrr}
\toprule
Estimator & Median ratio & Error & Within 10\% & Within 20\% & More than 30\% too large \\
\midrule
Object ratios, 35th percentile (current)          & 1.04 & 0.154 & 48\% & 74\% & 14\% \\
Object ratios, 30th percentile                    & 1.01 & 0.146 & 51\% & 76\% &  9\% \\
Object ratios, 25th percentile                    & 0.98 & 0.152 & 49\% & 76\% &  7\% \\
Object heights only, 35th percentile              & 1.06 & 0.202 & 48\% & 71\% & 18\% \\
Classes weighted by reliability, bias removed     & 0.95 & 0.155 & 45\% & 75\% &  8\% \\
Camera height only                                & 0.90 & 0.367 & 14\% & 40\% & 31\% \\
Object estimate and camera height, equal weights  & 0.96 & 0.216 & 28\% & 61\% & 20\% \\
\bottomrule
\end{tabular}
\end{table}

Should VGG-T-Det refine the scale? Today the detector cannot: it is given a ground-truth scale (the mean distance
of the scan points from the first camera) with its input, in training and at test time, so it never sees the
scale error. Two designs are possible. (i) The detector keeps receiving a scale, now the estimated one, and its
size regression has to absorb the error; this needs training with the estimated scale, or with the true scale
perturbed by the error distribution of Table~\ref{tab:sam3-q7}. (ii) The detector predicts one scale correction
per scene from all its queries together, which is applied to every box; a per-scene quantity is better
estimated from all objects at once than box by box. In both cases the question is whether the images and VGGT
tokens carry enough metric information beyond the object size priors the pipeline already uses.


\subsection{Answers so far to Q1, Q2 and Q5}
\label{sec:sam3-answers}

\paragraph{Q1: coverage of the detector's training data.}
Table~\ref{tab:sam3-a1} counts the annotated objects that fall into the 65 classes. A third (ScanNet) to more
than half (ScanNet++) of the annotated objects are outside them; these are mostly small clutter (book, paper,
box, bottle, socket, clothes, bag, light switch), spread over 3,126 distinct labels in ScanNet++. Only
ScanNet++ has boxes for all 65 classes. Seven classes exist in ScanNet++ alone: ceiling light (2,501 boxes),
power strip (241), mattress (138), cutting board (113), paper bag (110), coat rack (71) and shoe rack (49).
ARKitScenes has no boxes for 14 classes, among them counter, desk, monitor, closet and computer tower.
Whether SAM 3 finds the annotated objects in the training scenes is not measured yet.

\begin{table}[t]
\centering
\caption{Annotated objects of the three training datasets and the share that falls into the 65 classes.
Structure: walls, floors, ceilings and the generic ``object''.}
\label{tab:sam3-a1}
\begin{tabular}{lrrrrr}
\toprule
Dataset & Annotations & Structure & Objects & In the 65 classes & Classes with boxes \\
\midrule
ScanNet (EmbodiedScan, 1,201 scenes)     &  48,964 & 30.1\% & 34,220 & 66.4\% & 57 \\
ARKitScenes (EmbodiedScan, 4,498 scenes) & 102,623 & 16.3\% & 85,926 & 86.9\% & 51 \\
ScanNet++ (905 scenes)                   & 100,799 & 29.9\% & 70,672 & 42.4\% & 65 \\
\bottomrule
\end{tabular}
\end{table}

\paragraph{Q2: coverage of the Reason-190K questions.}
The objects a training question is about are the labels of its object list (Table~\ref{tab:sam3-a2};
room-size questions name none and are left out). Questions from ScanNet and ARKitScenes are covered almost
completely; the only notable gap is bicycle (287 ScanNet questions). Of the ScanNet++ questions, 28\% name an
object outside the 65 classes, most often light switch (2,337 questions), board (2,089), bottle, box, switch,
bag, slippers, book, blinds and jacket. That is about 18,400 of the 143,437 training questions (13\%).

\begin{table}[t]
\centering
\caption{Reason-190K questions whose named objects are all in the 65 classes. Exact: the class names and their
VSI-Bench synonyms. With renamings: also the dataset labels mapped to a class (ceiling lamp, shelf, mouse,
coat hanger, shoe, poster, office chair, couch).}
\label{tab:sam3-a2}
\begin{tabular}{lrrrr}
\toprule
Source & Questions & Exact & With renamings & Object mentions covered \\
\midrule
ScanNet     &  49,569 & 98.7\% & 98.9\% & 99.5\% \\
ARKitScenes &  28,485 & 99.1\% & 99.6\% & 99.8\% \\
ScanNet++   &  65,383 & 64.4\% & 71.8\% & 87.0\% \\
All         & 143,437 & 83.1\% & 86.7\% & 93.6\% \\
\bottomrule
\end{tabular}
\end{table}

\paragraph{Q5, wording: the proposed rewordings mostly do not help.}
Each rewording of Table~\ref{tab:sam3-q5-few} was run as a single SAM 3 prompt on the scenes whose questions
name the object and substituted for the benchmark name (Table~\ref{tab:sam3-a5}). Only one is a clear gain:
prompting ``radiator'' also where the question says ``heater'' raises the boxes from 6 to 21 for 46 objects and
the objects found from 11\% to 39\%. ``Ceiling lamp'', ``desktop computer'', ``PC case'' and ``extension cord with sockets'' find
fewer objects than the benchmark name. ``Kitchen counter'' and ``bag'' find nearly all objects (85\% and 88\%)
but add many extra boxes (14 for 13 counters, 22 for 43 backpacks), so they exchange one count error for the
other. ``Ventilation fan'' changes nothing. The benchmark names are thus mostly as good as or better than the alternatives, and the low recall of
ceiling light, computer tower, fan and power strip is not a wording problem.

\begin{table}[t]
\centering
\caption{Rewordings tested (box score $\geq 0.7$). Found: objects with a box of that name within 0.75\,m.
Exact: scenes with the right number of boxes.}
\label{tab:sam3-a5}
\begin{tabular}{llrrrrrr}
\toprule
Object (ground truth) & Prompt & Boxes & Extra & Missing & Found & Exact & Verdict \\
\midrule
heater / radiator (46) & benchmark name            &  6 &  0 & 40 & 11\% & 18\% & \\
                       & ``radiator''              & 21 &  1 & 26 & 39\% & 44\% & better \\
                       & ``wall radiator''         & 15 &  1 & 32 & 28\% & 35\% & better \\
ceiling light (116)    & benchmark name            & 53 &  1 & 64 & 39\% & 24\% & \\
                       & ``ceiling lamp''          & 21 &  0 & 95 & 16\% & 10\% & worse \\
computer tower (42)    & benchmark name            & 25 &  1 & 18 & 50\% & 14\% & \\
                       & ``desktop computer''      & 37 & 10 & 15 & 38\% & 14\% & worse \\
                       & ``PC case''               & 13 &  0 & 29 & 26\% &  0\% & worse \\
fan / exhaust fan (14) & benchmark name            &  5 &  1 & 10 & 21\% & 30\% & \\
                       & ``ventilation fan''       &  6 &  1 &  9 & 29\% & 30\% & no change \\
power strip (10)       & benchmark name            &  6 &  1 &  5 & 40\% & 29\% & \\
                       & ``extension cord with sockets'' & 0 & 0 & 10 & 0\% & 0\% & worse \\
counter (13)           & benchmark name            & 12 &  6 &  7 & 31\% & 18\% & \\
                       & ``kitchen counter''       & 27 & 14 &  0 & 85\% & 36\% & more found, more extra \\
backpack (43)          & benchmark name            & 29 &  0 & 14 & 67\% & 52\% & \\
                       & ``bag''                   & 64 & 22 &  1 & 88\% & 55\% & more found, more extra \\
                       & ``backpack'' and ``bag''  & 68 & 25 &  0 & 93\% & 52\% & more found, more extra \\
\bottomrule
\end{tabular}
\end{table}

\paragraph{Q5, decision.}
``Radiator'' is adopted: every scene whose questions say ``heater'' or ``radiator'' is prompted with ``radiator'',
and the box keeps the name the question uses. Over all 287 scenes this lowers the missing boxes of the list from
641 to 627 and leaves the extra boxes unchanged (414 to 415); the share of (scene, name) pairs with the exact
count rises from 54.2\% to 54.8\%. Radiators remain a weak class: 26 of 46 are still missing.
``Kitchen counter'' and ``bag'' are on hold: they find nearly all objects but add many extra boxes, so they are
useful only once a later stage removes extra boxes. The other rewordings are rejected. The prompt set is
therefore the names of the scene's questions with two substitutions, ``room door'' for door and ``radiator''
for heater.

\paragraph{Q5, extra boxes: half of them are on nothing that is annotated.}
Of 2,967 boxes with score $\geq 0.7$ in all scenes, 2,368 match a ground-truth object of their name and 599 do
not. Of these 599, 299 (50\%) lie where nothing is annotated, 163 (27\%) duplicate an object that already has a
box, and 137 (23\%) lie on an object of another class (Table~\ref{tab:sam3-a5x}). This corrects two guesses of
Table~\ref{tab:sam3-q5-many}. Duplicates are not the main cause for table, chair and sofa, so better merging
would remove about a quarter of the extra boxes at most; only for bed do duplicates lead. And tv boxes are
rarely on monitors (2 cases). Chair and sofa are confused in both directions, and pillow and cushion boxes often
land on the sofa or bed they lie on. Whether a box on nothing annotated is a real object without annotation or a
false detection cannot be told from the counts.

\begin{table}[t]
\centering
\caption{What the boxes of a name are (all 287 scenes, score $\geq 0.7$).}
\label{tab:sam3-a5x}
\begin{tabular}{lrrrrrl}
\toprule
Name & Boxes & Matched & Duplicate & Other class & Nothing annotated & Main other class \\
\midrule
table   & 454 & 317 & 48 & 25 & 64 & chair 11, bed 5 \\
chair   & 864 & 739 & 23 & 27 & 75 & sofa 17, table 6 \\
sofa    & 144 & 106 & 11 & 10 & 17 & chair 8 \\
door    & 163 & 131 &  6 &  2 & 24 & -- \\
bed     & 106 &  77 & 14 &  2 & 13 & -- \\
pillow  &  93 &  70 &  9 & 13 &  1 & sofa 6, bed 4 \\
tv      &  96 &  78 &  8 &  5 &  5 & monitor 2 \\
cushion &  12 &   5 &  1 &  6 &  0 & sofa 3, chair 3 \\
\midrule
All names & 2,967 & 2,368 & 163 & 137 & 299 & \\
\bottomrule
\end{tabular}
\end{table}

\subsection{True scale against estimated scale}
\label{sec:sam3-scale}

The SAM 3 + VGGT-SLAM list estimates its metric scale from object sizes (Q7). To measure what that estimate costs, the
same list (prompts from the scene's questions, ``room door'', ``radiator'') was scored once as it is and once with
its scale replaced by the true one, obtained by fitting the reconstructed cameras to the ground-truth cameras.
Nothing else changes: the same boxes, only rescaled about the first camera. The estimated scale is a median 10\%
off (49\% of the scenes within 10\%, 87\% within 30\%).

\paragraph{3D boxes.} Table~\ref{tab:sam3-scale-3d}: predicted boxes are matched to ground-truth boxes of the same
label (287 VSI-Bench scenes, 3,194 objects); precision counts boxes with score $\geq 0.7$. The true scale improves
every loose-match measure by 5 to 9 points, on all three datasets (Table~\ref{tab:sam3-scale-ds}). Tight matches do
not improve (mAP@0.50 stays near 4): the scale only fixes the distance of every box from the camera, not the boxes'
shape and placement errors.

\begin{table}[t]
\centering
\caption{SAM 3 + VGGT-SLAM boxes with the estimated and the true scale (287 VSI-Bench scenes, 3,194 objects).}
\label{tab:sam3-scale-3d}
\begin{tabular}{lrrr}
\toprule
 & Estimated scale & True scale & Change \\
\midrule
Recall, IoU $\geq 0.25$             & 30.5 & 38.2 & $+7.7$ \\
Precision, IoU $\geq 0.25$          & 32.5 & 40.3 & $+7.8$ \\
mAP@0.25                            & 17.2 & 22.2 & $+5.0$ \\
Recall, IoU $\geq 0.5$              &  8.1 & 11.2 & $+3.1$ \\
mAP@0.50                            &  4.4 &  4.0 & $-0.4$ \\
Recall, centre within 0.5\,m        & 67.3 & 75.9 & $+8.6$ \\
Precision, centre within 0.5\,m     & 65.8 & 73.9 & $+8.1$ \\
\bottomrule
\end{tabular}
\end{table}

\begin{table}[t]
\centering
\caption{Recall at IoU $\geq 0.25$ by dataset.}
\label{tab:sam3-scale-ds}
\begin{tabular}{lrrr}
\toprule
Dataset & Objects & Estimated scale & True scale \\
\midrule
ScanNet     & 1,224 & 32.0 & 41.9 \\
ARKitScenes &   854 & 34.8 & 43.0 \\
ScanNet++   & 1,116 & 25.4 & 30.4 \\
\bottomrule
\end{tabular}
\end{table}

\paragraph{VSI-Bench.} Table~\ref{tab:sam3-scale-vsi} gives the score of the default language model (mix-v2-1) with
the two lists. The true scale helps the two metric tasks a little (absolute distance $+1.0$, size $+1.3$) and leaves
the overall score unchanged; counting even drops by 1.1, within the noise of a single run. The scale error therefore
costs far less than the errors in the list itself: which objects are listed and how loose their boxes are.

\begin{table}[t]
\centering
\caption{VSI-Bench (all 5,130 questions) with the SAM 3 + VGGT-SLAM list at the estimated and the true scale.}
\label{tab:sam3-scale-vsi}
\begin{tabular}{lrrrrr}
\toprule
SAM 3 list & Full & Debiased & Abs.\ distance & Size & Counting \\
\midrule
Estimated scale & 69.0 & 62.9 & 57.5 & 73.4 & 73.5 \\
True scale      & 69.1 & 63.0 & 58.5 & 74.7 & 72.4 \\
\bottomrule
\end{tabular}
\end{table}

\paragraph{Which results use the true scale.} Every VGG-T-Det variant (the released 18-class model, cls-31, cls-63,
cls-65 and IndoorDet-65-31k) multiplies VGGT's normalized point cloud by the mean distance of the ground-truth scan
from the first camera and uses the first camera's ground-truth pose, in training and at test time; the proposals of
IndoorDet-65-31k are placed with the true scale as well. All object lists built from a detector, and all language-model
training data (ground-truth and detector boxes), are therefore metric with the true scale. Only the SAM 3 + VGGT-SLAM
lists use an estimated scale, so they are the only video-only lists, and the language model sees a scale error only
at test time with them.

Table~\ref{tab:sam3-scale-use} lists which runs use the true scale.

\begin{table}[t]
\centering
\caption{Runs with and without the true scale. ``True scale'' = metric scale from the ground-truth scan (detectors) or
from fitting the reconstructed cameras to the ground-truth cameras (SAM 3 lists); ``estimated'' = from object size
priors, with up direction and walls from the reconstruction and the first camera as origin (no ground truth).}
\label{tab:sam3-scale-use}
\small
\begin{tabular}{llll}
\toprule
Run & Stage & Scale & Used for \\
\midrule
\multicolumn{4}{l}{\emph{Without the true scale (video only)}} \\
SAM 3 + VGGT-SLAM lists, all prompt variants (fixed 65 names, question & test & estimated & VSI-Bench, e.g.\ 69.0 / 62.9 \\
\quad names, ``room door'', ``radiator'', route phrases, thresholds 0.1--0.9) & & & \\
VGGT-SLAM lists with oracle 2D boxes (\texttt{slam\_manh}, \texttt{slamall\_manh}) & test & estimated & VSI-Bench \\
SAM 3 list ablation: predicted centre and size cells & test & estimated & VSI-Bench (Section~\ref{sec:sam3}) \\
SAM 3 list, estimated-scale arm of Table~\ref{tab:sam3-scale-3d} & test & estimated & 3D evaluation, VSI-Bench \\
Video-only model & train, test & no boxes & VSI-Bench 69.4 / 63.7 \\
\midrule
\multicolumn{4}{l}{\emph{With the true scale}} \\
VGG-T-Det: released 18-class, ES-56, cls-31, VSI-40, cls-63, cls-65 & train, test & true & detection; VSI-Bench lists \\
IndoorDet-65-31k (test2-hybrid), and its SAM 3 proposals & train, test & true & detection; VSI-Bench lists \\
Lists from a detector (cls-31; IndoorDet-65-31k refine and & test & true & VSI-Bench, e.g.\ 69.9 / 62.7, 69.8 / 63.2 \\
\quad threshold lists) & & & \\
VGGT-SLAM lists \texttt{\_truescale}; SAM 3 list, true-scale arm & test & true & VSI-Bench, 3D evaluation \\
Language-model SFT data (mix-v2-0, mix-v2-1, Exp1): ground-truth & train & true & every SFT run \\
\quad and detector boxes & & & \\
\bottomrule
\end{tabular}
\end{table}
; needs \usepackage{booktabs}.
% All numbers: VSI-Bench, 5,130 questions, 288 scenes (287 processed); box score threshold 0.7 unless stated.
% Sources: VGGT-SLAM/exps/vsibench_objects (s3_alt_stats.py, s3_badlist_thr.py, qvocab.py, s3_rows.py),
% 3d-scene-lang/runs/eval_q35vs_gandalf/{mixv2_1_sam3_full,mixv2_1_s3q_roomdoor}_rvs_question.

\section{Open-vocabulary object lists with SAM 3}
\label{sec:sam3}

\paragraph{Pipeline.}
SAM 3 is run on up to 100 VGGT-SLAM keyframes per video with one text prompt per object name
(per-frame confidence $\geq 0.5$). Each mask is lifted to 3D with the SLAM points, masks of the same name are
merged across frames, and each merged instance gives one box. The box score is the mean of the three best
SAM 3 scores times $1-\exp(-\mathrm{views}/2)$. No ground truth is used. The resulting list of
\{label, box, score\} entries is given to the language model together with the video.

\paragraph{The object list, not the geometry, limits the benefit.}
Replacing the predicted centres and sizes by ground truth while keeping the SAM 3 list raises the VSI-Bench
score only from 68.6 to 69.9, whereas the full ground-truth list reaches 77.3. The errors that matter are
extra boxes and missing objects. With question-name prompts and ``room door'' (below), the list has 414 extra
and 641 missing boxes against 3,194 ground-truth objects of the names the questions use.

\subsection{Run time}
\label{sec:sam3-time}

Table~\ref{tab:sam3-time} gives the cost of the two stages, measured on 597 ScanNet++ scenes processed together
on four NVIDIA RTX PRO 6000 Blackwell Max-Q Workstation Edition GPUs (96\,GB each, machine gandalf) with three
workers per GPU; the GPUs ran nothing else. VGGT-SLAM receives up to 600 frames per scene and
keeps 100 keyframes; SAM 3 runs all 65 class prompts on each of the 100 keyframes (per-frame confidence 0.3).
SAM 3 costs about eight times as much as VGGT-SLAM per scene. Its time hardly varies between scenes, since every
prompt is run on every keyframe (about 4\,s per keyframe for 65 prompts); with three workers sharing a GPU it is
limited by GPU compute, so more workers per GPU do not raise the throughput. For comparison, SAM 3 on ARKitScenes on
four NVIDIA RTX 6000 Ada GPUs (46\,GB each, machine wukong; two workers per GPU) took about
250\,s per scene per worker with 51 prompts.

\begin{table}[t]
\centering
\caption{Run time per scene of the two stages (597 ScanNet++ scenes, 4 NVIDIA RTX PRO 6000 Blackwell Max-Q GPUs,
96\,GB, 12 workers).}
\label{tab:sam3-time}
\begin{tabular}{lrrrr}
\toprule
Stage & \multicolumn{3}{c}{Time per scene per worker} & Throughput \\
      & Median & Mean & 90th percentile & (12 workers) \\
\midrule
VGGT-SLAM (600 frames in, 100 keyframes)   &  48\,s &  48\,s &  74\,s & 10 scenes per minute \\
SAM 3 (65 prompts on 100 keyframes)        & 394\,s & 392\,s & 398\,s & 110 scenes per hour \\
\bottomrule
\end{tabular}
\end{table}

\subsection{Prompt wording changes the number of boxes}

\paragraph{``door'' versus ``room door''.}
The prompt ``door'' also fires on cabinet, wardrobe and shower doors. It alone produced 37\% of all extra boxes
(231 of 620). Prompting ``room door'' instead removes most of them at the cost of some recall
(Table~\ref{tab:sam3-door}).

\begin{table}[t]
\centering
\caption{Door boxes for the two prompts (170 ground-truth doors in 114 scenes).
Found: a box of that label within 0.75\,m of the ground-truth centre.}
\label{tab:sam3-door}
\begin{tabular}{lrrrrr}
\toprule
Prompt & Boxes & Extra & Missing & Scenes with exact count & Found \\
\midrule
``door''      & 391 & 231 & 10 & 31\% & 85\% \\
``room door'' & 163 &  25 & 32 & 64\% & 76\% \\
\bottomrule
\end{tabular}
\end{table}

\paragraph{A more specific word is not always better.}
For tables, every rewording we tried is worse than the plain name (Table~\ref{tab:sam3-table}): ``desk'' removes
13 extra boxes but loses 61 more tables. The extra table boxes are duplicates of real tables and tables without
annotation, so they call for better merging, not for another word.

\begin{table}[t]
\centering
\caption{Table boxes for four prompts (406 ground-truth tables, 285 scenes, all with ``room door'').}
\label{tab:sam3-table}
\begin{tabular}{lrrrrr}
\toprule
Prompt & Boxes & Extra & Missing & Scenes with exact count & Found \\
\midrule
``table'' (benchmark name)      & 450 & 101 &  57 & 41\% & 75\% \\
``desk''                        & 376 &  88 & 118 & 35\% & 62\% \\
``table or desk''               & 529 & 162 &  39 & 34\% & 79\% \\
``dining table or office desk'' & 287 &  58 & 177 & 36\% & 50\% \\
\bottomrule
\end{tabular}
\end{table}

\paragraph{Which names are detected too rarely or too often.}
Table~\ref{tab:sam3-names} lists the names with the largest count errors. Two groups appear.
Names detected too rarely are mostly ones whose benchmark wording is unusual for the object
(``heater'' for a wall radiator, ``ceiling light'', ``computer tower'') or small objects (basket, shoes, bucket,
pan, kettle), which have no extra boxes at all. Names detected too often (table, chair, sofa, bed) are large
furniture that is split into several boxes or duplicated. ``pillow'' and ``cushion'' fire on the same objects.
Only the door and table wordings have been tested; the others are candidates.

\begin{table}[t]
\centering
\caption{Count errors per object name (question-name prompts with ``room door''). Scenes: those whose questions
use the name and that contain the object.}
\label{tab:sam3-names}
\begin{tabular}{lrrrrr}
\toprule
Name & Ground truth & Boxes & Extra & Missing & Scenes with too few \\
\midrule
\multicolumn{6}{l}{\emph{Too few boxes}} \\
heater / radiator & 46 & 6 & 0 & 40 & 82\% \\
ceiling light     & 116 & 53 & 1 & 64 & 72\% \\
computer tower    & 42 & 25 & 1 & 18 & 71\% \\
fan / exhaust fan & 14 & 5 & 1 & 10 & 60\% \\
power strip       & 10 & 6 & 1 & 5 & 57\% \\
backpack          & 43 & 29 & 0 & 14 & 48\% \\
window            & 83 & 68 & 11 & 26 & 40\% \\
\midrule
\multicolumn{6}{l}{\emph{Too many boxes}} \\
table & 409 & 454 & 102 & 57 & 23\% \\
chair & 897 & 864 & 73 & 106 & 26\% \\
sofa  & 146 & 144 & 30 & 32 & 31\% \\
bed   & 82 & 106 & 26 & 2 & 3\% \\
door (``room door'') & 170 & 163 & 25 & 32 & 18\% \\
\bottomrule
\end{tabular}
\end{table}

\paragraph{The score threshold trades missing objects for extra ones.}
Lowering the box-score threshold does not fix the list (Table~\ref{tab:sam3-thr}, fixed 65-name prompts): at 0.3
fewer labels are short of boxes, but there is about one extra box per ground-truth object. About 7\% of labels
are never found at any threshold, because SAM 3 gives no mask above its per-frame confidence.

\begin{table}[t]
\centering
\caption{Per (scene, label) count against ground truth for four score thresholds.}
\label{tab:sam3-thr}
\begin{tabular}{rrrrr}
\toprule
Threshold & Exact & Too few & Too many & Extra boxes per ground-truth object \\
\midrule
0.7 & 44.4\% & 22.6\% & 33.0\% & 0.43 \\
0.5 & 41.6\% & 15.6\% & 42.8\% & 0.64 \\
0.3 & 30.9\% &  9.9\% & 59.2\% & 1.06 \\
0.2 & 24.2\% &  7.0\% & 68.8\% & 1.48 \\
\bottomrule
\end{tabular}
\end{table}

\subsection{The vocabulary of the questions}

\paragraph{Seven of eight tasks use a closed list of 65 names.}
Outside route planning, every object in a VSI-Bench question is one of 65 names, and the wording depends on the
source dataset (``trash can'' / ``trash bin'', ``heater'' / ``radiator'', ``whiteboard'' / ``blackboard'').
Prompting SAM 3 with the names a scene's own questions use, in their original wording, gives 15.9 boxes per
question instead of 41.9 for a fixed list of all names.

\paragraph{Route-planning questions are open-vocabulary.}
The 194 route-planning questions are free text. Of their 979 object mentions (start, facing direction, goal and
every ``go forward until''), only 505 (52\%) are one of the 65 names. The rest are other objects
(sink in 28 scenes, cabinet 11, bathroom sink 8, desk 6, doorframe 6, wardrobe 4),
objects with an attribute (``blue chair'', ``open door'', ``glass coffee table''),
spelling variants (``trashbin'', ``bathtube'', ``refridgerator'', ``dinning table'')
and places that are not objects (wall, ``into the room'', ``outside'').
A fixed vocabulary therefore gives no box at all for about half of the objects these questions refer to.

\paragraph{Extraction rule.}
We take each mention in the question's own words, remove articles, position words and trailing relations
(``the sink on the left'' $\to$ ``sink'', ``foot of the bed'' $\to$ ``bed''), keep attributes
(``blue chair''), and drop walls, rooms and ``outside''. This gives 183 distinct phrases; 241 (scene, phrase)
pairs in 147 scenes are not already question names and become additional prompts. Only a plain ``door'' is
prompted as ``room door''; ``doorframe'', ``doorway'', ``open door'' or ``kitchen door'' are used as written.
When a route phrase and a benchmark name find the same object, the list keeps one box, shown under the name the
question uses.

\subsection{Effect on VSI-Bench}

Table~\ref{tab:sam3-vsi} compares the lists with the same language model (Qwen3.5-4B, mix-v2-1).
Question-name prompts with ``room door'' gain 0.4 / 0.5 points over the fixed 65-name prompts. The gain is in
counting ($+1.5$) and absolute distance ($+0.8$), the tasks most affected by extra door boxes; the other tasks
change by at most 0.6. The list still does not beat the video alone.
Adding the route phrases changes the boxes shown for 130 of the 194 route-planning questions, yet the
route-planning score stays at 52.1: only 10 answers change, 4 from wrong to right and 4 from right to wrong.
Section~\ref{sec:sam3-route} examines why.

\begin{table}[t]
\centering
\caption{VSI-Bench score (mean over tasks) on all 5,130 questions and on the debiased subset.}
\label{tab:sam3-vsi}
\begin{tabular}{lrrrrr}
\toprule
Object list & Full & Debiased & Counting & Abs.\ distance & Route planning \\
\midrule
None (video only)                               & 69.4 & 63.7 & -- & -- & -- \\
SAM 3, fixed 65 names                           & 68.6 & 62.3 & 72.0 & 56.6 & 51.5 \\
SAM 3, question names + ``room door''           & 69.0 & 62.8 & 73.5 & 57.4 & 52.1 \\
SAM 3, question names + ``room door'' + route phrases & 69.0 & 62.7 & 73.3 & 57.4 & 52.1 \\
SAM 3 list with ground-truth centre and size    & 69.9 & 63.9 & -- & -- & -- \\
Ground-truth list                               & 77.3 & 73.9 & -- & -- & -- \\
\bottomrule
\end{tabular}
\end{table}

\subsection{Route planning: are the named objects in the list?}
\label{sec:sam3-route}

Route planning is the task an open vocabulary should help most, since half of its object mentions are outside
the 65 names. We therefore check, for each of the 194 questions, whether every object it names has a box in
the list that includes the route phrases. A question names 659 distinct objects in total, 418 by a benchmark
name and 241 by another phrase.

\begin{table}[t]
\centering
\caption{Route-planning questions: objects named in the question that have at least one box of that name,
for four box-score thresholds. The evaluated list uses 0.7.}
\label{tab:sam3-route}
\begin{tabular}{rrrr}
\toprule
Threshold & Questions with every object found & Benchmark names found & Other phrases found \\
\midrule
0.7  & 120 / 194 (62\%) & 375 / 418 (90\%) & 198 / 241 (82\%) \\
0.5  & 145 / 194 (75\%) & 394 / 418 (94\%) & 210 / 241 (87\%) \\
0.3  & 157 / 194 (81\%) & 400 / 418 (96\%) & 218 / 241 (90\%) \\
none & 164 / 194 (85\%) & 403 / 418 (96\%) & 222 / 241 (92\%) \\
\bottomrule
\end{tabular}
\end{table}

\paragraph{Coverage is incomplete.}
At the evaluated threshold only 62\% of the questions have a box for every object they name
(Table~\ref{tab:sam3-route}). The objects missing most often are heater (9 questions), door (7), window (4),
cabinet (3) and shower (3). Heater is also the main object never found at any threshold (8 questions),
in line with Table~\ref{tab:sam3-names}; the remaining never-found objects are single cases such as
``column'', ``trophies'' or ``paper shredder box''.

\paragraph{Found objects are often ambiguous.}
A route question refers to one specific object (``the door'', ``the blue chair''). At threshold 0.7,
149 of the 375 found benchmark names and 91 of the 198 found phrases have several boxes of that name, so the
list does not say which box is meant. Lowering the threshold closes part of the coverage gap but widens this
one: with no threshold, 295 of the 403 found benchmark names have several boxes.

\paragraph{The language model barely uses the added boxes.}
The route phrases change the boxes shown for 130 of the 194 questions, and the 6 questions that matched no box
before now match some. Yet only 10 answers change, 4 from wrong to right and 4 from right to wrong; 101 answers
are correct with and without the phrases (52.1\%). A likely reason, which we have not
tested, is that the model was fine-tuned on lists that use the dataset label names, so boxes under free phrases
are an input it has not been trained to use.

\paragraph{Limits of this check.}
``Found'' means that a box with the name exists. Free phrases have no ground-truth boxes, so we cannot verify
that the box is on the right object. Mentions that are not objects (walls, rooms, ``outside'') are not counted.

\paragraph{Conclusion for route planning.}
Three gaps act together: 38\% of the questions miss at least one named object, about 40\% of the found objects
are ambiguous, and the model does not exploit the additional boxes. Extending the vocabulary alone therefore
gives no gain (52.1 with and without route phrases); the task needs boxes that identify the referred instance
and a model trained on such lists.

\paragraph{Summary.}
(i) Text prompts are not interchangeable: one word (``room'') removes 206 of 231 extra door boxes, while a
plausible synonym for table loses 61 objects.
(ii) Benchmark wording should be kept per dataset, and ambiguous names disambiguated, not replaced by synonyms.
(iii) A closed vocabulary covers only about half of the objects in free-text questions; taking the prompts from
the question raises coverage to 82--90\% of the named objects, but does not change the route-planning score
(Section~\ref{sec:sam3-route}).
(iv) Wording fixes a minority of the errors: after them, missing objects (641) outnumber extra boxes (414), and
neither a lower threshold nor rewording removes both.

\subsection{Open questions}
\label{sec:sam3-open}

\paragraph{Q1. Do the 65 classes (with ``room door'') cover the detector's training data?}
The detector (VGG-T-Det) is trained on ScanNet, ScanNet++ and ARKitScenes. Do the 65 class names, prompted as
in Section~\ref{sec:sam3} with ``room door'' for door, cover the annotated objects of these three datasets, and
which annotated objects fall outside them?
\emph{Known so far:} the 65 classes are the VSI-Bench names, so every class has a benchmark use, but the
converse is not measured. In the 242 ScanNet++ scenes that are ready, 7,881 of 27,878 annotations map to one of
the 65 classes (the rest include walls, floors and the generic ``object''); seven classes have boxes only in
ScanNet++ (ceiling light, cutting board, paper bag, shoe rack, coat rack, power strip, mattress) and two have
none there (fireplace, piano). The shares for ScanNet and ARKitScenes, and whether SAM 3 with these prompts
finds the annotated objects of each class in the training scenes, are still to be measured.

\paragraph{Q2. Do the 65 classes cover the objects named in the Reason-190K questions?}
The language model is fine-tuned on Reason-190K. If its questions name objects outside the 65 classes, a
65-class list gives no box for them in training, as it does for route planning in VSI-Bench.
\emph{Known so far:} not measured for the 65-class set. The measurement is the one of
Section~\ref{sec:sam3-route} applied to the training questions: the share of questions whose named objects are
all among the 65 names, per source dataset.

\paragraph{Q3. Should the whole pipeline be open-vocabulary?}
If the answers to Q1 and Q2 show large gaps, the alternative is to drop the fixed class set: the label of a box
comes from the SAM 3 prompt, the classification head of the detector is removed, and the detector only refines
the centre and size of each box and decides whether to keep it. For this to work the detector must receive
more candidate boxes than are finally kept, so the SAM 3 score threshold would be lowered and the removal of
extra boxes left to the keep decision.
\emph{Known so far:} lowering the box-score threshold from 0.7 to 0.3 reduces the labels with too few boxes
from 22.6\% to 9.9\% but raises the extra boxes per ground-truth object from 0.43 to 1.06
(Table~\ref{tab:sam3-thr}), so a keep decision would have to remove about one box per object. About 7\% of
labels are found at no box-score threshold; recovering them needs a lower SAM 3 per-frame confidence (0.5 now),
which has not been tried. Two risks are open: a correct box on an object without annotation is trained as
``drop'', and the detector cannot add an object that SAM 3 missed.

\paragraph{Q4. Should SAM 3 run on overlapping video segments that are merged afterwards, like submaps?}
VGGT-SLAM handles long videos by reconstructing submaps of consecutive frames that share overlapping frames and
then aligning them. The same scheme could be used for the objects: run SAM 3 on each video segment, with the
overlapping frames in common, build the object instances per segment, and merge the segment instances at the end,
using the overlapping frames to decide which instances are the same object.
\emph{Known so far:} the current pipeline does not use segments. SAM 3 is applied as an image model to each
keyframe independently (at most 100 keyframes spread evenly over the video, so long videos are subsampled), and
masks are associated only afterwards in 3D, by the overlap of their lifted points over the whole video. Extra
boxes from this global merge are the main count error for large furniture (Table~\ref{tab:sam3-names}), and
errors of the global reconstruction enter every association. Open points: whether SAM 3's video mode, which
tracks an object through consecutive frames, gives cleaner instances within a segment than per-frame masks;
whether merging per segment and then across segments produces fewer duplicates than one global merge; and how
many keyframes per segment are needed so that small objects are seen in several views.

\paragraph{Q5. Which count errors can be fixed by wording, and which need better merging?}
Two groups of names have wrong box counts (question-name prompts with ``room door'', box score $\geq 0.7$).
For the names with too few boxes, the benchmark wording is the likely cause, and
Table~\ref{tab:sam3-q5-few} lists a wording to try for each. For the names with too many boxes, the likely
cause is the merging of masks into instances, not the wording (Table~\ref{tab:sam3-q5-many}).
\emph{Known so far:} only the door and table wordings have been tested
(Tables~\ref{tab:sam3-door} and~\ref{tab:sam3-table}). The wordings and causes in the two tables below are
hypotheses read from the counts; none has been tested. Small objects with too few boxes and no extra boxes
(basket, shoes, bucket, pan, kettle, blanket, bowl, cup) are probably not a wording problem; a lower SAM 3
per-frame confidence is the more likely remedy.

\begin{table}[t]
\centering
\caption{Names with too few boxes, where the wording is the likely cause, and the wording to try (untested).}
\label{tab:sam3-q5-few}
\begin{tabular}{lrrrl}
\toprule
Name & Ground truth & Boxes & Scenes with too few & Wording to try \\
\midrule
heater / radiator & 46  & 6  & 82\% & ``radiator'' for both, or ``wall radiator'' \\
ceiling light     & 116 & 53 & 72\% & ``ceiling lamp'' (the ScanNet++ annotators' word) \\
computer tower    & 42  & 25 & 71\% & ``desktop computer'', ``PC case'' \\
fan / exhaust fan & 14  & 5  & 60\% & ``ventilation fan'' \\
power strip       & 10  & 6  & 57\% & ``extension cord with sockets'' \\
counter           & 13  & 12 & 55\% & ``kitchen counter'' \\
backpack          & 43  & 29 & 48\% & ``bag'' as a second prompt \\
window            & 83  & 68 & 40\% & unclear; may be visibility, not wording \\
\bottomrule
\end{tabular}
\end{table}

\begin{table}[t]
\centering
\caption{Names with too many boxes, where merging rather than wording is the likely cause (untested).}
\label{tab:sam3-q5-many}
\begin{tabular}{lrrrl}
\toprule
Name & Ground truth & Boxes & Extra & Likely cause \\
\midrule
table            & 409 & 454 & 102 & duplicates, tables without annotation \\
chair            & 897 & 864 & 73  & duplicates (also 106 missing) \\
sofa             & 146 & 144 & 30  & sectional sofas split into pieces \\
bed              & 82  & 106 & 26  & bed split into several boxes \\
tv               & 94  & 96  & 15  & monitors detected as tv \\
pillow / cushion & 100 & 105 & 20  & the two names detect the same objects \\
\bottomrule
\end{tabular}
\end{table}

\paragraph{Q6. Which box-score threshold is best for the language model?}
All VSI-Bench results above use boxes with score $\geq 0.7$. The count statistics
(Table~\ref{tab:sam3-thr}) show that a lower threshold trades missing objects for extra boxes, but not which side
of the trade the language model prefers: every box is given with its score and the model is told to trust a box
according to it, so it may tolerate low-score boxes. We evaluate the same list (prompts = the names in the
scene's questions, ``room door'' for door) at thresholds 0.1, 0.3, 0.5, 0.7 and 0.9 on all 5,130 questions with
the same model.
\emph{Result:} the threshold hardly matters (Table~\ref{tab:sam3-q6}). Over a threefold change in the
number of boxes, the full score stays between 68.3 and 69.2 and the debiased score between 62.2 and 62.8; 0.7 is
best or within 0.2 of the best on both. The one clear trend is counting, which falls steadily as the threshold
is lowered (74.0 at 0.9, 69.7 at 0.1): extra boxes are counted. Relative distance (66.3 to 68.9), absolute distance and
route planning change without a trend. No threshold reaches the video-only score (69.4 / 63.7), so the
threshold is not the lever; lowering it is useful only if a later stage removes the extra boxes.

\begin{table}[t]
\centering
\caption{VSI-Bench score for five box-score thresholds (question names + ``room door''). Boxes: mean number of
boxes in the list of a scene before the list is restricted to the objects a question names.}
\label{tab:sam3-q6}
\begin{tabular}{rrrrrrrr}
\toprule
Threshold & Boxes & Full & Debiased & Counting & Abs.\ distance & Rel.\ distance & Route planning \\
\midrule
0.1 & 32.6 & 68.3 & 62.2 & 69.7 & 56.8 & 68.0 & 52.1 \\
0.3 & 26.4 & 68.7 & 62.7 & 70.9 & 56.8 & 67.7 & 53.1 \\
0.5 & 19.5 & 68.5 & 62.2 & 72.2 & 57.3 & 66.3 & 50.5 \\
0.7 & 15.9 & 69.0 & 62.8 & 73.5 & 57.4 & 66.9 & 52.1 \\
0.9 & 10.0 & 69.2 & 62.2 & 74.0 & 56.7 & 68.9 & 52.1 \\
\bottomrule
\end{tabular}
\end{table}

\paragraph{Q7. The estimated scale is too large on average. Can it be improved, and should VGG-T-Det refine it?}
VGGT-SLAM reconstructs a scene up to an unknown scale. The pipeline estimates the metric scale from the objects:
for every well-seen instance of a class with a size prior, the ratio of the prior footprint and height to the
measured ones is taken, and the scale is the 35th percentile of these ratios (a low percentile, because partly
seen objects measure too small and would otherwise inflate the scale).
\emph{Known so far:} against the true scale (from the ground-truth camera poses), the estimate is too large in
the median by 4\% and in the mean by 11\% (Table~\ref{tab:sam3-q7}). The error is skewed: a tenth of the scenes
is more than a third too large, a tenth more than 13\% too small. Only 48\% of the scenes are within 10\% and
74\% within 20\%. The scale error is the largest geometric error of the list: in the ablation of
Section~\ref{sec:sam3}, ground-truth box sizes are the only geometric correction that helps the language model,
and in the query check for the detector the estimated scale costs 9 to 14 points of box coverage compared with
the true scale.

\begin{table}[t]
\centering
\caption{Estimated scale divided by true scale, per scene (question names + ``room door'').}
\label{tab:sam3-q7}
\begin{tabular}{lrrrrrr}
\toprule
Dataset & Scenes & Median & Mean & 10th--90th percentile & Within 10\% & Within 20\% \\
\midrule
ScanNet     &  88 & 1.06 & 1.14 & 0.89--1.33 & 49\% & 78\% \\
ARKitScenes & 150 & 1.05 & 1.10 & 0.86--1.35 & 43\% & 70\% \\
ScanNet++   &  49 & 1.03 & 1.07 & 0.83--1.30 & 65\% & 78\% \\
All         & 287 & 1.04 & 1.11 & 0.87--1.34 & 48\% & 74\% \\
\bottomrule
\end{tabular}
\end{table}

Ways to improve the estimate (the first three are tested below):
\begin{itemize}
\item \emph{Camera height.} The height of the camera above the floor is measured in the reconstruction and is
      close to a constant for hand-held video; it gives a second, object-free scale estimate. It can be combined
      with the object ratios or used to reject outliers.
\item \emph{More reliable objects.} Use only instances seen completely (mask not cut by the image border, many
      views), prefer classes with a narrow size distribution (door height, toilet, bed length) over classes with
      a wide one (table, cabinet), and weight the ratios by the spread of the class prior.
\item \emph{Calibrate the percentile.} The 35th percentile was set by hand; the bias in Table~\ref{tab:sam3-q7}
      suggests a lower one, and it can be fitted on the training scenes, per dataset if needed.
\item \emph{A learned metric prior.} A monocular metric-depth model on a few keyframes gives the scale without
      object priors.
\end{itemize}

\emph{Tested without training} (Table~\ref{tab:sam3-q7b}; constants fitted on 939 ScanNet training scenes, evaluated on
the 287 VSI-Bench scenes; the ground truth is used for the evaluation only).
The bias is easy to remove: taking the 30th instead of the 35th percentile of the object ratios centres the
estimate (median 1.01) and halves the share of scenes that are more than 30\% too large (14\% to 9\%). The spread
hardly changes: the typical error stays near 15\% and only half of the scenes are within 10\%. The other ideas
do not help. The camera height is not a usable constant: its true value is 0.8\,m in ScanNet and 1.3\,m in
ARKitScenes and ScanNet++, and it varies from 0.5 to 1.4\,m within ScanNet, so a scale from the camera height
alone is far worse, and mixing it with the object estimate is worse than the object estimate alone. Weighting
the classes by how reliable their size is, with the per-class bias removed, gives the same error as the plain
percentile. The limit is therefore the size priors themselves: object sizes vary too much within a class to fix
the scale to better than about 15\% from a handful of objects. A better scale needs another source of metric
information, a metric-depth model or the detector.

\begin{table}[t]
\centering
\caption{Scale estimators that use no ground truth. Error: mean absolute log ratio of estimated to true scale.}
\label{tab:sam3-q7b}
\begin{tabular}{lrrrrr}
\toprule
Estimator & Median ratio & Error & Within 10\% & Within 20\% & More than 30\% too large \\
\midrule
Object ratios, 35th percentile (current)          & 1.04 & 0.154 & 48\% & 74\% & 14\% \\
Object ratios, 30th percentile                    & 1.01 & 0.146 & 51\% & 76\% &  9\% \\
Object ratios, 25th percentile                    & 0.98 & 0.152 & 49\% & 76\% &  7\% \\
Object heights only, 35th percentile              & 1.06 & 0.202 & 48\% & 71\% & 18\% \\
Classes weighted by reliability, bias removed     & 0.95 & 0.155 & 45\% & 75\% &  8\% \\
Camera height only                                & 0.90 & 0.367 & 14\% & 40\% & 31\% \\
Object estimate and camera height, equal weights  & 0.96 & 0.216 & 28\% & 61\% & 20\% \\
\bottomrule
\end{tabular}
\end{table}

Should VGG-T-Det refine the scale? Today the detector cannot: it is given a ground-truth scale (the mean distance
of the scan points from the first camera) with its input, in training and at test time, so it never sees the
scale error. Two designs are possible. (i) The detector keeps receiving a scale, now the estimated one, and its
size regression has to absorb the error; this needs training with the estimated scale, or with the true scale
perturbed by the error distribution of Table~\ref{tab:sam3-q7}. (ii) The detector predicts one scale correction
per scene from all its queries together, which is applied to every box; a per-scene quantity is better
estimated from all objects at once than box by box. In both cases the question is whether the images and VGGT
tokens carry enough metric information beyond the object size priors the pipeline already uses.


\subsection{Answers so far to Q1, Q2 and Q5}
\label{sec:sam3-answers}

\paragraph{Q1: coverage of the detector's training data.}
Table~\ref{tab:sam3-a1} counts the annotated objects that fall into the 65 classes. A third (ScanNet) to more
than half (ScanNet++) of the annotated objects are outside them; these are mostly small clutter (book, paper,
box, bottle, socket, clothes, bag, light switch), spread over 3,126 distinct labels in ScanNet++. Only
ScanNet++ has boxes for all 65 classes. Seven classes exist in ScanNet++ alone: ceiling light (2,501 boxes),
power strip (241), mattress (138), cutting board (113), paper bag (110), coat rack (71) and shoe rack (49).
ARKitScenes has no boxes for 14 classes, among them counter, desk, monitor, closet and computer tower.
Whether SAM 3 finds the annotated objects in the training scenes is not measured yet.

\begin{table}[t]
\centering
\caption{Annotated objects of the three training datasets and the share that falls into the 65 classes.
Structure: walls, floors, ceilings and the generic ``object''.}
\label{tab:sam3-a1}
\begin{tabular}{lrrrrr}
\toprule
Dataset & Annotations & Structure & Objects & In the 65 classes & Classes with boxes \\
\midrule
ScanNet (EmbodiedScan, 1,201 scenes)     &  48,964 & 30.1\% & 34,220 & 66.4\% & 57 \\
ARKitScenes (EmbodiedScan, 4,498 scenes) & 102,623 & 16.3\% & 85,926 & 86.9\% & 51 \\
ScanNet++ (905 scenes)                   & 100,799 & 29.9\% & 70,672 & 42.4\% & 65 \\
\bottomrule
\end{tabular}
\end{table}

\paragraph{Q2: coverage of the Reason-190K questions.}
The objects a training question is about are the labels of its object list (Table~\ref{tab:sam3-a2};
room-size questions name none and are left out). Questions from ScanNet and ARKitScenes are covered almost
completely; the only notable gap is bicycle (287 ScanNet questions). Of the ScanNet++ questions, 28\% name an
object outside the 65 classes, most often light switch (2,337 questions), board (2,089), bottle, box, switch,
bag, slippers, book, blinds and jacket. That is about 18,400 of the 143,437 training questions (13\%).

\begin{table}[t]
\centering
\caption{Reason-190K questions whose named objects are all in the 65 classes. Exact: the class names and their
VSI-Bench synonyms. With renamings: also the dataset labels mapped to a class (ceiling lamp, shelf, mouse,
coat hanger, shoe, poster, office chair, couch).}
\label{tab:sam3-a2}
\begin{tabular}{lrrrr}
\toprule
Source & Questions & Exact & With renamings & Object mentions covered \\
\midrule
ScanNet     &  49,569 & 98.7\% & 98.9\% & 99.5\% \\
ARKitScenes &  28,485 & 99.1\% & 99.6\% & 99.8\% \\
ScanNet++   &  65,383 & 64.4\% & 71.8\% & 87.0\% \\
All         & 143,437 & 83.1\% & 86.7\% & 93.6\% \\
\bottomrule
\end{tabular}
\end{table}

\paragraph{Q5, wording: the proposed rewordings mostly do not help.}
Each rewording of Table~\ref{tab:sam3-q5-few} was run as a single SAM 3 prompt on the scenes whose questions
name the object and substituted for the benchmark name (Table~\ref{tab:sam3-a5}). Only one is a clear gain:
prompting ``radiator'' also where the question says ``heater'' raises the boxes from 6 to 21 for 46 objects and
the objects found from 11\% to 39\%. ``Ceiling lamp'', ``desktop computer'', ``PC case'' and ``extension cord with sockets'' find
fewer objects than the benchmark name. ``Kitchen counter'' and ``bag'' find nearly all objects (85\% and 88\%)
but add many extra boxes (14 for 13 counters, 22 for 43 backpacks), so they exchange one count error for the
other. ``Ventilation fan'' changes nothing. The benchmark names are thus mostly as good as or better than the alternatives, and the low recall of
ceiling light, computer tower, fan and power strip is not a wording problem.

\begin{table}[t]
\centering
\caption{Rewordings tested (box score $\geq 0.7$). Found: objects with a box of that name within 0.75\,m.
Exact: scenes with the right number of boxes.}
\label{tab:sam3-a5}
\begin{tabular}{llrrrrrr}
\toprule
Object (ground truth) & Prompt & Boxes & Extra & Missing & Found & Exact & Verdict \\
\midrule
heater / radiator (46) & benchmark name            &  6 &  0 & 40 & 11\% & 18\% & \\
                       & ``radiator''              & 21 &  1 & 26 & 39\% & 44\% & better \\
                       & ``wall radiator''         & 15 &  1 & 32 & 28\% & 35\% & better \\
ceiling light (116)    & benchmark name            & 53 &  1 & 64 & 39\% & 24\% & \\
                       & ``ceiling lamp''          & 21 &  0 & 95 & 16\% & 10\% & worse \\
computer tower (42)    & benchmark name            & 25 &  1 & 18 & 50\% & 14\% & \\
                       & ``desktop computer''      & 37 & 10 & 15 & 38\% & 14\% & worse \\
                       & ``PC case''               & 13 &  0 & 29 & 26\% &  0\% & worse \\
fan / exhaust fan (14) & benchmark name            &  5 &  1 & 10 & 21\% & 30\% & \\
                       & ``ventilation fan''       &  6 &  1 &  9 & 29\% & 30\% & no change \\
power strip (10)       & benchmark name            &  6 &  1 &  5 & 40\% & 29\% & \\
                       & ``extension cord with sockets'' & 0 & 0 & 10 & 0\% & 0\% & worse \\
counter (13)           & benchmark name            & 12 &  6 &  7 & 31\% & 18\% & \\
                       & ``kitchen counter''       & 27 & 14 &  0 & 85\% & 36\% & more found, more extra \\
backpack (43)          & benchmark name            & 29 &  0 & 14 & 67\% & 52\% & \\
                       & ``bag''                   & 64 & 22 &  1 & 88\% & 55\% & more found, more extra \\
                       & ``backpack'' and ``bag''  & 68 & 25 &  0 & 93\% & 52\% & more found, more extra \\
\bottomrule
\end{tabular}
\end{table}

\paragraph{Q5, decision.}
``Radiator'' is adopted: every scene whose questions say ``heater'' or ``radiator'' is prompted with ``radiator'',
and the box keeps the name the question uses. Over all 287 scenes this lowers the missing boxes of the list from
641 to 627 and leaves the extra boxes unchanged (414 to 415); the share of (scene, name) pairs with the exact
count rises from 54.2\% to 54.8\%. Radiators remain a weak class: 26 of 46 are still missing.
``Kitchen counter'' and ``bag'' are on hold: they find nearly all objects but add many extra boxes, so they are
useful only once a later stage removes extra boxes. The other rewordings are rejected. The prompt set is
therefore the names of the scene's questions with two substitutions, ``room door'' for door and ``radiator''
for heater.

\paragraph{Q5, extra boxes: half of them are on nothing that is annotated.}
Of 2,967 boxes with score $\geq 0.7$ in all scenes, 2,368 match a ground-truth object of their name and 599 do
not. Of these 599, 299 (50\%) lie where nothing is annotated, 163 (27\%) duplicate an object that already has a
box, and 137 (23\%) lie on an object of another class (Table~\ref{tab:sam3-a5x}). This corrects two guesses of
Table~\ref{tab:sam3-q5-many}. Duplicates are not the main cause for table, chair and sofa, so better merging
would remove about a quarter of the extra boxes at most; only for bed do duplicates lead. And tv boxes are
rarely on monitors (2 cases). Chair and sofa are confused in both directions, and pillow and cushion boxes often
land on the sofa or bed they lie on. Whether a box on nothing annotated is a real object without annotation or a
false detection cannot be told from the counts.

\begin{table}[t]
\centering
\caption{What the boxes of a name are (all 287 scenes, score $\geq 0.7$).}
\label{tab:sam3-a5x}
\begin{tabular}{lrrrrrl}
\toprule
Name & Boxes & Matched & Duplicate & Other class & Nothing annotated & Main other class \\
\midrule
table   & 454 & 317 & 48 & 25 & 64 & chair 11, bed 5 \\
chair   & 864 & 739 & 23 & 27 & 75 & sofa 17, table 6 \\
sofa    & 144 & 106 & 11 & 10 & 17 & chair 8 \\
door    & 163 & 131 &  6 &  2 & 24 & -- \\
bed     & 106 &  77 & 14 &  2 & 13 & -- \\
pillow  &  93 &  70 &  9 & 13 &  1 & sofa 6, bed 4 \\
tv      &  96 &  78 &  8 &  5 &  5 & monitor 2 \\
cushion &  12 &   5 &  1 &  6 &  0 & sofa 3, chair 3 \\
\midrule
All names & 2,967 & 2,368 & 163 & 137 & 299 & \\
\bottomrule
\end{tabular}
\end{table}

\subsection{True scale against estimated scale}
\label{sec:sam3-scale}

The SAM 3 + VGGT-SLAM list estimates its metric scale from object sizes (Q7). To measure what that estimate costs, the
same list (prompts from the scene's questions, ``room door'', ``radiator'') was scored once as it is and once with
its scale replaced by the true one, obtained by fitting the reconstructed cameras to the ground-truth cameras.
Nothing else changes: the same boxes, only rescaled about the first camera. The estimated scale is a median 10\%
off (49\% of the scenes within 10\%, 87\% within 30\%).

\paragraph{3D boxes.} Table~\ref{tab:sam3-scale-3d}: predicted boxes are matched to ground-truth boxes of the same
label (287 VSI-Bench scenes, 3,194 objects); precision counts boxes with score $\geq 0.7$. The true scale improves
every loose-match measure by 5 to 9 points, on all three datasets (Table~\ref{tab:sam3-scale-ds}). Tight matches do
not improve (mAP@0.50 stays near 4): the scale only fixes the distance of every box from the camera, not the boxes'
shape and placement errors.

\begin{table}[t]
\centering
\caption{SAM 3 + VGGT-SLAM boxes with the estimated and the true scale (287 VSI-Bench scenes, 3,194 objects).}
\label{tab:sam3-scale-3d}
\begin{tabular}{lrrr}
\toprule
 & Estimated scale & True scale & Change \\
\midrule
Recall, IoU $\geq 0.25$             & 30.5 & 38.2 & $+7.7$ \\
Precision, IoU $\geq 0.25$          & 32.5 & 40.3 & $+7.8$ \\
mAP@0.25                            & 17.2 & 22.2 & $+5.0$ \\
Recall, IoU $\geq 0.5$              &  8.1 & 11.2 & $+3.1$ \\
mAP@0.50                            &  4.4 &  4.0 & $-0.4$ \\
Recall, centre within 0.5\,m        & 67.3 & 75.9 & $+8.6$ \\
Precision, centre within 0.5\,m     & 65.8 & 73.9 & $+8.1$ \\
\bottomrule
\end{tabular}
\end{table}

\begin{table}[t]
\centering
\caption{Recall at IoU $\geq 0.25$ by dataset.}
\label{tab:sam3-scale-ds}
\begin{tabular}{lrrr}
\toprule
Dataset & Objects & Estimated scale & True scale \\
\midrule
ScanNet     & 1,224 & 32.0 & 41.9 \\
ARKitScenes &   854 & 34.8 & 43.0 \\
ScanNet++   & 1,116 & 25.4 & 30.4 \\
\bottomrule
\end{tabular}
\end{table}

\paragraph{VSI-Bench.} Table~\ref{tab:sam3-scale-vsi} gives the score of the default language model (mix-v2-1) with
the two lists. The true scale helps the two metric tasks a little (absolute distance $+1.0$, size $+1.3$) and leaves
the overall score unchanged; counting even drops by 1.1, within the noise of a single run. The scale error therefore
costs far less than the errors in the list itself: which objects are listed and how loose their boxes are.

\begin{table}[t]
\centering
\caption{VSI-Bench (all 5,130 questions) with the SAM 3 + VGGT-SLAM list at the estimated and the true scale.}
\label{tab:sam3-scale-vsi}
\begin{tabular}{lrrrrr}
\toprule
SAM 3 list & Full & Debiased & Abs.\ distance & Size & Counting \\
\midrule
Estimated scale & 69.0 & 62.9 & 57.5 & 73.4 & 73.5 \\
True scale      & 69.1 & 63.0 & 58.5 & 74.7 & 72.4 \\
\bottomrule
\end{tabular}
\end{table}

\paragraph{Which results use the true scale.} Every VGG-T-Det variant (the released 18-class model, cls-31, cls-63,
cls-65 and IndoorDet-65-31k) multiplies VGGT's normalized point cloud by the mean distance of the ground-truth scan
from the first camera and uses the first camera's ground-truth pose, in training and at test time; the proposals of
IndoorDet-65-31k are placed with the true scale as well. All object lists built from a detector, and all language-model
training data (ground-truth and detector boxes), are therefore metric with the true scale. Only the SAM 3 + VGGT-SLAM
lists use an estimated scale, so they are the only video-only lists, and the language model sees a scale error only
at test time with them.

Table~\ref{tab:sam3-scale-use} lists which runs use the true scale.

\begin{table}[t]
\centering
\caption{Runs with and without the true scale. ``True scale'' = metric scale from the ground-truth scan (detectors) or
from fitting the reconstructed cameras to the ground-truth cameras (SAM 3 lists); ``estimated'' = from object size
priors, with up direction and walls from the reconstruction and the first camera as origin (no ground truth).}
\label{tab:sam3-scale-use}
\small
\begin{tabular}{llll}
\toprule
Run & Stage & Scale & Used for \\
\midrule
\multicolumn{4}{l}{\emph{Without the true scale (video only)}} \\
SAM 3 + VGGT-SLAM lists, all prompt variants (fixed 65 names, question & test & estimated & VSI-Bench, e.g.\ 69.0 / 62.9 \\
\quad names, ``room door'', ``radiator'', route phrases, thresholds 0.1--0.9) & & & \\
VGGT-SLAM lists with oracle 2D boxes (\texttt{slam\_manh}, \texttt{slamall\_manh}) & test & estimated & VSI-Bench \\
SAM 3 list ablation: predicted centre and size cells & test & estimated & VSI-Bench (Section~\ref{sec:sam3}) \\
SAM 3 list, estimated-scale arm of Table~\ref{tab:sam3-scale-3d} & test & estimated & 3D evaluation, VSI-Bench \\
Video-only model & train, test & no boxes & VSI-Bench 69.4 / 63.7 \\
\midrule
\multicolumn{4}{l}{\emph{With the true scale}} \\
VGG-T-Det: released 18-class, ES-56, cls-31, VSI-40, cls-63, cls-65 & train, test & true & detection; VSI-Bench lists \\
IndoorDet-65-31k (test2-hybrid), and its SAM 3 proposals & train, test & true & detection; VSI-Bench lists \\
Lists from a detector (cls-31; IndoorDet-65-31k refine and & test & true & VSI-Bench, e.g.\ 69.9 / 62.7, 69.8 / 63.2 \\
\quad threshold lists) & & & \\
VGGT-SLAM lists \texttt{\_truescale}; SAM 3 list, true-scale arm & test & true & VSI-Bench, 3D evaluation \\
Language-model SFT data (mix-v2-0, mix-v2-1, Exp1): ground-truth & train & true & every SFT run \\
\quad and detector boxes & & & \\
\bottomrule
\end{tabular}
\end{table}
; needs \usepackage{booktabs}.
% All numbers: VSI-Bench, 5,130 questions, 288 scenes (287 processed); box score threshold 0.7 unless stated.
% Sources: VGGT-SLAM/exps/vsibench_objects (s3_alt_stats.py, s3_badlist_thr.py, qvocab.py, s3_rows.py),
% 3d-scene-lang/runs/eval_q35vs_gandalf/{mixv2_1_sam3_full,mixv2_1_s3q_roomdoor}_rvs_question.

\section{Open-vocabulary object lists with SAM 3}
\label{sec:sam3}

\paragraph{Pipeline.}
SAM 3 is run on up to 100 VGGT-SLAM keyframes per video with one text prompt per object name
(per-frame confidence $\geq 0.5$). Each mask is lifted to 3D with the SLAM points, masks of the same name are
merged across frames, and each merged instance gives one box. The box score is the mean of the three best
SAM 3 scores times $1-\exp(-\mathrm{views}/2)$. No ground truth is used. The resulting list of
\{label, box, score\} entries is given to the language model together with the video.

\paragraph{The object list, not the geometry, limits the benefit.}
Replacing the predicted centres and sizes by ground truth while keeping the SAM 3 list raises the VSI-Bench
score only from 68.6 to 69.9, whereas the full ground-truth list reaches 77.3. The errors that matter are
extra boxes and missing objects. With question-name prompts and ``room door'' (below), the list has 414 extra
and 641 missing boxes against 3,194 ground-truth objects of the names the questions use.

\subsection{Run time}
\label{sec:sam3-time}

Table~\ref{tab:sam3-time} gives the cost of the two stages, measured on 597 ScanNet++ scenes processed together
on four NVIDIA RTX PRO 6000 Blackwell Max-Q Workstation Edition GPUs (96\,GB each, machine gandalf) with three
workers per GPU; the GPUs ran nothing else. VGGT-SLAM receives up to 600 frames per scene and
keeps 100 keyframes; SAM 3 runs all 65 class prompts on each of the 100 keyframes (per-frame confidence 0.3).
SAM 3 costs about eight times as much as VGGT-SLAM per scene. Its time hardly varies between scenes, since every
prompt is run on every keyframe (about 4\,s per keyframe for 65 prompts); with three workers sharing a GPU it is
limited by GPU compute, so more workers per GPU do not raise the throughput. For comparison, SAM 3 on ARKitScenes on
four NVIDIA RTX 6000 Ada GPUs (46\,GB each, machine wukong; two workers per GPU) took about
250\,s per scene per worker with 51 prompts.

\begin{table}[t]
\centering
\caption{Run time per scene of the two stages (597 ScanNet++ scenes, 4 NVIDIA RTX PRO 6000 Blackwell Max-Q GPUs,
96\,GB, 12 workers).}
\label{tab:sam3-time}
\begin{tabular}{lrrrr}
\toprule
Stage & \multicolumn{3}{c}{Time per scene per worker} & Throughput \\
      & Median & Mean & 90th percentile & (12 workers) \\
\midrule
VGGT-SLAM (600 frames in, 100 keyframes)   &  48\,s &  48\,s &  74\,s & 10 scenes per minute \\
SAM 3 (65 prompts on 100 keyframes)        & 394\,s & 392\,s & 398\,s & 110 scenes per hour \\
\bottomrule
\end{tabular}
\end{table}

\subsection{Prompt wording changes the number of boxes}

\paragraph{``door'' versus ``room door''.}
The prompt ``door'' also fires on cabinet, wardrobe and shower doors. It alone produced 37\% of all extra boxes
(231 of 620). Prompting ``room door'' instead removes most of them at the cost of some recall
(Table~\ref{tab:sam3-door}).

\begin{table}[t]
\centering
\caption{Door boxes for the two prompts (170 ground-truth doors in 114 scenes).
Found: a box of that label within 0.75\,m of the ground-truth centre.}
\label{tab:sam3-door}
\begin{tabular}{lrrrrr}
\toprule
Prompt & Boxes & Extra & Missing & Scenes with exact count & Found \\
\midrule
``door''      & 391 & 231 & 10 & 31\% & 85\% \\
``room door'' & 163 &  25 & 32 & 64\% & 76\% \\
\bottomrule
\end{tabular}
\end{table}

\paragraph{A more specific word is not always better.}
For tables, every rewording we tried is worse than the plain name (Table~\ref{tab:sam3-table}): ``desk'' removes
13 extra boxes but loses 61 more tables. The extra table boxes are duplicates of real tables and tables without
annotation, so they call for better merging, not for another word.

\begin{table}[t]
\centering
\caption{Table boxes for four prompts (406 ground-truth tables, 285 scenes, all with ``room door'').}
\label{tab:sam3-table}
\begin{tabular}{lrrrrr}
\toprule
Prompt & Boxes & Extra & Missing & Scenes with exact count & Found \\
\midrule
``table'' (benchmark name)      & 450 & 101 &  57 & 41\% & 75\% \\
``desk''                        & 376 &  88 & 118 & 35\% & 62\% \\
``table or desk''               & 529 & 162 &  39 & 34\% & 79\% \\
``dining table or office desk'' & 287 &  58 & 177 & 36\% & 50\% \\
\bottomrule
\end{tabular}
\end{table}

\paragraph{Which names are detected too rarely or too often.}
Table~\ref{tab:sam3-names} lists the names with the largest count errors. Two groups appear.
Names detected too rarely are mostly ones whose benchmark wording is unusual for the object
(``heater'' for a wall radiator, ``ceiling light'', ``computer tower'') or small objects (basket, shoes, bucket,
pan, kettle), which have no extra boxes at all. Names detected too often (table, chair, sofa, bed) are large
furniture that is split into several boxes or duplicated. ``pillow'' and ``cushion'' fire on the same objects.
Only the door and table wordings have been tested; the others are candidates.

\begin{table}[t]
\centering
\caption{Count errors per object name (question-name prompts with ``room door''). Scenes: those whose questions
use the name and that contain the object.}
\label{tab:sam3-names}
\begin{tabular}{lrrrrr}
\toprule
Name & Ground truth & Boxes & Extra & Missing & Scenes with too few \\
\midrule
\multicolumn{6}{l}{\emph{Too few boxes}} \\
heater / radiator & 46 & 6 & 0 & 40 & 82\% \\
ceiling light     & 116 & 53 & 1 & 64 & 72\% \\
computer tower    & 42 & 25 & 1 & 18 & 71\% \\
fan / exhaust fan & 14 & 5 & 1 & 10 & 60\% \\
power strip       & 10 & 6 & 1 & 5 & 57\% \\
backpack          & 43 & 29 & 0 & 14 & 48\% \\
window            & 83 & 68 & 11 & 26 & 40\% \\
\midrule
\multicolumn{6}{l}{\emph{Too many boxes}} \\
table & 409 & 454 & 102 & 57 & 23\% \\
chair & 897 & 864 & 73 & 106 & 26\% \\
sofa  & 146 & 144 & 30 & 32 & 31\% \\
bed   & 82 & 106 & 26 & 2 & 3\% \\
door (``room door'') & 170 & 163 & 25 & 32 & 18\% \\
\bottomrule
\end{tabular}
\end{table}

\paragraph{The score threshold trades missing objects for extra ones.}
Lowering the box-score threshold does not fix the list (Table~\ref{tab:sam3-thr}, fixed 65-name prompts): at 0.3
fewer labels are short of boxes, but there is about one extra box per ground-truth object. About 7\% of labels
are never found at any threshold, because SAM 3 gives no mask above its per-frame confidence.

\begin{table}[t]
\centering
\caption{Per (scene, label) count against ground truth for four score thresholds.}
\label{tab:sam3-thr}
\begin{tabular}{rrrrr}
\toprule
Threshold & Exact & Too few & Too many & Extra boxes per ground-truth object \\
\midrule
0.7 & 44.4\% & 22.6\% & 33.0\% & 0.43 \\
0.5 & 41.6\% & 15.6\% & 42.8\% & 0.64 \\
0.3 & 30.9\% &  9.9\% & 59.2\% & 1.06 \\
0.2 & 24.2\% &  7.0\% & 68.8\% & 1.48 \\
\bottomrule
\end{tabular}
\end{table}

\subsection{The vocabulary of the questions}

\paragraph{Seven of eight tasks use a closed list of 65 names.}
Outside route planning, every object in a VSI-Bench question is one of 65 names, and the wording depends on the
source dataset (``trash can'' / ``trash bin'', ``heater'' / ``radiator'', ``whiteboard'' / ``blackboard'').
Prompting SAM 3 with the names a scene's own questions use, in their original wording, gives 15.9 boxes per
question instead of 41.9 for a fixed list of all names.

\paragraph{Route-planning questions are open-vocabulary.}
The 194 route-planning questions are free text. Of their 979 object mentions (start, facing direction, goal and
every ``go forward until''), only 505 (52\%) are one of the 65 names. The rest are other objects
(sink in 28 scenes, cabinet 11, bathroom sink 8, desk 6, doorframe 6, wardrobe 4),
objects with an attribute (``blue chair'', ``open door'', ``glass coffee table''),
spelling variants (``trashbin'', ``bathtube'', ``refridgerator'', ``dinning table'')
and places that are not objects (wall, ``into the room'', ``outside'').
A fixed vocabulary therefore gives no box at all for about half of the objects these questions refer to.

\paragraph{Extraction rule.}
We take each mention in the question's own words, remove articles, position words and trailing relations
(``the sink on the left'' $\to$ ``sink'', ``foot of the bed'' $\to$ ``bed''), keep attributes
(``blue chair''), and drop walls, rooms and ``outside''. This gives 183 distinct phrases; 241 (scene, phrase)
pairs in 147 scenes are not already question names and become additional prompts. Only a plain ``door'' is
prompted as ``room door''; ``doorframe'', ``doorway'', ``open door'' or ``kitchen door'' are used as written.
When a route phrase and a benchmark name find the same object, the list keeps one box, shown under the name the
question uses.

\subsection{Effect on VSI-Bench}

Table~\ref{tab:sam3-vsi} compares the lists with the same language model (Qwen3.5-4B, mix-v2-1).
Question-name prompts with ``room door'' gain 0.4 / 0.5 points over the fixed 65-name prompts. The gain is in
counting ($+1.5$) and absolute distance ($+0.8$), the tasks most affected by extra door boxes; the other tasks
change by at most 0.6. The list still does not beat the video alone.
Adding the route phrases changes the boxes shown for 130 of the 194 route-planning questions, yet the
route-planning score stays at 52.1: only 10 answers change, 4 from wrong to right and 4 from right to wrong.
Section~\ref{sec:sam3-route} examines why.

\begin{table}[t]
\centering
\caption{VSI-Bench score (mean over tasks) on all 5,130 questions and on the debiased subset.}
\label{tab:sam3-vsi}
\begin{tabular}{lrrrrr}
\toprule
Object list & Full & Debiased & Counting & Abs.\ distance & Route planning \\
\midrule
None (video only)                               & 69.4 & 63.7 & -- & -- & -- \\
SAM 3, fixed 65 names                           & 68.6 & 62.3 & 72.0 & 56.6 & 51.5 \\
SAM 3, question names + ``room door''           & 69.0 & 62.8 & 73.5 & 57.4 & 52.1 \\
SAM 3, question names + ``room door'' + route phrases & 69.0 & 62.7 & 73.3 & 57.4 & 52.1 \\
SAM 3 list with ground-truth centre and size    & 69.9 & 63.9 & -- & -- & -- \\
Ground-truth list                               & 77.3 & 73.9 & -- & -- & -- \\
\bottomrule
\end{tabular}
\end{table}

\subsection{Route planning: are the named objects in the list?}
\label{sec:sam3-route}

Route planning is the task an open vocabulary should help most, since half of its object mentions are outside
the 65 names. We therefore check, for each of the 194 questions, whether every object it names has a box in
the list that includes the route phrases. A question names 659 distinct objects in total, 418 by a benchmark
name and 241 by another phrase.

\begin{table}[t]
\centering
\caption{Route-planning questions: objects named in the question that have at least one box of that name,
for four box-score thresholds. The evaluated list uses 0.7.}
\label{tab:sam3-route}
\begin{tabular}{rrrr}
\toprule
Threshold & Questions with every object found & Benchmark names found & Other phrases found \\
\midrule
0.7  & 120 / 194 (62\%) & 375 / 418 (90\%) & 198 / 241 (82\%) \\
0.5  & 145 / 194 (75\%) & 394 / 418 (94\%) & 210 / 241 (87\%) \\
0.3  & 157 / 194 (81\%) & 400 / 418 (96\%) & 218 / 241 (90\%) \\
none & 164 / 194 (85\%) & 403 / 418 (96\%) & 222 / 241 (92\%) \\
\bottomrule
\end{tabular}
\end{table}

\paragraph{Coverage is incomplete.}
At the evaluated threshold only 62\% of the questions have a box for every object they name
(Table~\ref{tab:sam3-route}). The objects missing most often are heater (9 questions), door (7), window (4),
cabinet (3) and shower (3). Heater is also the main object never found at any threshold (8 questions),
in line with Table~\ref{tab:sam3-names}; the remaining never-found objects are single cases such as
``column'', ``trophies'' or ``paper shredder box''.

\paragraph{Found objects are often ambiguous.}
A route question refers to one specific object (``the door'', ``the blue chair''). At threshold 0.7,
149 of the 375 found benchmark names and 91 of the 198 found phrases have several boxes of that name, so the
list does not say which box is meant. Lowering the threshold closes part of the coverage gap but widens this
one: with no threshold, 295 of the 403 found benchmark names have several boxes.

\paragraph{The language model barely uses the added boxes.}
The route phrases change the boxes shown for 130 of the 194 questions, and the 6 questions that matched no box
before now match some. Yet only 10 answers change, 4 from wrong to right and 4 from right to wrong; 101 answers
are correct with and without the phrases (52.1\%). A likely reason, which we have not
tested, is that the model was fine-tuned on lists that use the dataset label names, so boxes under free phrases
are an input it has not been trained to use.

\paragraph{Limits of this check.}
``Found'' means that a box with the name exists. Free phrases have no ground-truth boxes, so we cannot verify
that the box is on the right object. Mentions that are not objects (walls, rooms, ``outside'') are not counted.

\paragraph{Conclusion for route planning.}
Three gaps act together: 38\% of the questions miss at least one named object, about 40\% of the found objects
are ambiguous, and the model does not exploit the additional boxes. Extending the vocabulary alone therefore
gives no gain (52.1 with and without route phrases); the task needs boxes that identify the referred instance
and a model trained on such lists.

\paragraph{Summary.}
(i) Text prompts are not interchangeable: one word (``room'') removes 206 of 231 extra door boxes, while a
plausible synonym for table loses 61 objects.
(ii) Benchmark wording should be kept per dataset, and ambiguous names disambiguated, not replaced by synonyms.
(iii) A closed vocabulary covers only about half of the objects in free-text questions; taking the prompts from
the question raises coverage to 82--90\% of the named objects, but does not change the route-planning score
(Section~\ref{sec:sam3-route}).
(iv) Wording fixes a minority of the errors: after them, missing objects (641) outnumber extra boxes (414), and
neither a lower threshold nor rewording removes both.

\subsection{Open questions}
\label{sec:sam3-open}

\paragraph{Q1. Do the 65 classes (with ``room door'') cover the detector's training data?}
The detector (VGG-T-Det) is trained on ScanNet, ScanNet++ and ARKitScenes. Do the 65 class names, prompted as
in Section~\ref{sec:sam3} with ``room door'' for door, cover the annotated objects of these three datasets, and
which annotated objects fall outside them?
\emph{Known so far:} the 65 classes are the VSI-Bench names, so every class has a benchmark use, but the
converse is not measured. In the 242 ScanNet++ scenes that are ready, 7,881 of 27,878 annotations map to one of
the 65 classes (the rest include walls, floors and the generic ``object''); seven classes have boxes only in
ScanNet++ (ceiling light, cutting board, paper bag, shoe rack, coat rack, power strip, mattress) and two have
none there (fireplace, piano). The shares for ScanNet and ARKitScenes, and whether SAM 3 with these prompts
finds the annotated objects of each class in the training scenes, are still to be measured.

\paragraph{Q2. Do the 65 classes cover the objects named in the Reason-190K questions?}
The language model is fine-tuned on Reason-190K. If its questions name objects outside the 65 classes, a
65-class list gives no box for them in training, as it does for route planning in VSI-Bench.
\emph{Known so far:} not measured for the 65-class set. The measurement is the one of
Section~\ref{sec:sam3-route} applied to the training questions: the share of questions whose named objects are
all among the 65 names, per source dataset.

\paragraph{Q3. Should the whole pipeline be open-vocabulary?}
If the answers to Q1 and Q2 show large gaps, the alternative is to drop the fixed class set: the label of a box
comes from the SAM 3 prompt, the classification head of the detector is removed, and the detector only refines
the centre and size of each box and decides whether to keep it. For this to work the detector must receive
more candidate boxes than are finally kept, so the SAM 3 score threshold would be lowered and the removal of
extra boxes left to the keep decision.
\emph{Known so far:} lowering the box-score threshold from 0.7 to 0.3 reduces the labels with too few boxes
from 22.6\% to 9.9\% but raises the extra boxes per ground-truth object from 0.43 to 1.06
(Table~\ref{tab:sam3-thr}), so a keep decision would have to remove about one box per object. About 7\% of
labels are found at no box-score threshold; recovering them needs a lower SAM 3 per-frame confidence (0.5 now),
which has not been tried. Two risks are open: a correct box on an object without annotation is trained as
``drop'', and the detector cannot add an object that SAM 3 missed.

\paragraph{Q4. Should SAM 3 run on overlapping video segments that are merged afterwards, like submaps?}
VGGT-SLAM handles long videos by reconstructing submaps of consecutive frames that share overlapping frames and
then aligning them. The same scheme could be used for the objects: run SAM 3 on each video segment, with the
overlapping frames in common, build the object instances per segment, and merge the segment instances at the end,
using the overlapping frames to decide which instances are the same object.
\emph{Known so far:} the current pipeline does not use segments. SAM 3 is applied as an image model to each
keyframe independently (at most 100 keyframes spread evenly over the video, so long videos are subsampled), and
masks are associated only afterwards in 3D, by the overlap of their lifted points over the whole video. Extra
boxes from this global merge are the main count error for large furniture (Table~\ref{tab:sam3-names}), and
errors of the global reconstruction enter every association. Open points: whether SAM 3's video mode, which
tracks an object through consecutive frames, gives cleaner instances within a segment than per-frame masks;
whether merging per segment and then across segments produces fewer duplicates than one global merge; and how
many keyframes per segment are needed so that small objects are seen in several views.

\paragraph{Q5. Which count errors can be fixed by wording, and which need better merging?}
Two groups of names have wrong box counts (question-name prompts with ``room door'', box score $\geq 0.7$).
For the names with too few boxes, the benchmark wording is the likely cause, and
Table~\ref{tab:sam3-q5-few} lists a wording to try for each. For the names with too many boxes, the likely
cause is the merging of masks into instances, not the wording (Table~\ref{tab:sam3-q5-many}).
\emph{Known so far:} only the door and table wordings have been tested
(Tables~\ref{tab:sam3-door} and~\ref{tab:sam3-table}). The wordings and causes in the two tables below are
hypotheses read from the counts; none has been tested. Small objects with too few boxes and no extra boxes
(basket, shoes, bucket, pan, kettle, blanket, bowl, cup) are probably not a wording problem; a lower SAM 3
per-frame confidence is the more likely remedy.

\begin{table}[t]
\centering
\caption{Names with too few boxes, where the wording is the likely cause, and the wording to try (untested).}
\label{tab:sam3-q5-few}
\begin{tabular}{lrrrl}
\toprule
Name & Ground truth & Boxes & Scenes with too few & Wording to try \\
\midrule
heater / radiator & 46  & 6  & 82\% & ``radiator'' for both, or ``wall radiator'' \\
ceiling light     & 116 & 53 & 72\% & ``ceiling lamp'' (the ScanNet++ annotators' word) \\
computer tower    & 42  & 25 & 71\% & ``desktop computer'', ``PC case'' \\
fan / exhaust fan & 14  & 5  & 60\% & ``ventilation fan'' \\
power strip       & 10  & 6  & 57\% & ``extension cord with sockets'' \\
counter           & 13  & 12 & 55\% & ``kitchen counter'' \\
backpack          & 43  & 29 & 48\% & ``bag'' as a second prompt \\
window            & 83  & 68 & 40\% & unclear; may be visibility, not wording \\
\bottomrule
\end{tabular}
\end{table}

\begin{table}[t]
\centering
\caption{Names with too many boxes, where merging rather than wording is the likely cause (untested).}
\label{tab:sam3-q5-many}
\begin{tabular}{lrrrl}
\toprule
Name & Ground truth & Boxes & Extra & Likely cause \\
\midrule
table            & 409 & 454 & 102 & duplicates, tables without annotation \\
chair            & 897 & 864 & 73  & duplicates (also 106 missing) \\
sofa             & 146 & 144 & 30  & sectional sofas split into pieces \\
bed              & 82  & 106 & 26  & bed split into several boxes \\
tv               & 94  & 96  & 15  & monitors detected as tv \\
pillow / cushion & 100 & 105 & 20  & the two names detect the same objects \\
\bottomrule
\end{tabular}
\end{table}

\paragraph{Q6. Which box-score threshold is best for the language model?}
All VSI-Bench results above use boxes with score $\geq 0.7$. The count statistics
(Table~\ref{tab:sam3-thr}) show that a lower threshold trades missing objects for extra boxes, but not which side
of the trade the language model prefers: every box is given with its score and the model is told to trust a box
according to it, so it may tolerate low-score boxes. We evaluate the same list (prompts = the names in the
scene's questions, ``room door'' for door) at thresholds 0.1, 0.3, 0.5, 0.7 and 0.9 on all 5,130 questions with
the same model.
\emph{Result:} the threshold hardly matters (Table~\ref{tab:sam3-q6}). Over a threefold change in the
number of boxes, the full score stays between 68.3 and 69.2 and the debiased score between 62.2 and 62.8; 0.7 is
best or within 0.2 of the best on both. The one clear trend is counting, which falls steadily as the threshold
is lowered (74.0 at 0.9, 69.7 at 0.1): extra boxes are counted. Relative distance (66.3 to 68.9), absolute distance and
route planning change without a trend. No threshold reaches the video-only score (69.4 / 63.7), so the
threshold is not the lever; lowering it is useful only if a later stage removes the extra boxes.

\begin{table}[t]
\centering
\caption{VSI-Bench score for five box-score thresholds (question names + ``room door''). Boxes: mean number of
boxes in the list of a scene before the list is restricted to the objects a question names.}
\label{tab:sam3-q6}
\begin{tabular}{rrrrrrrr}
\toprule
Threshold & Boxes & Full & Debiased & Counting & Abs.\ distance & Rel.\ distance & Route planning \\
\midrule
0.1 & 32.6 & 68.3 & 62.2 & 69.7 & 56.8 & 68.0 & 52.1 \\
0.3 & 26.4 & 68.7 & 62.7 & 70.9 & 56.8 & 67.7 & 53.1 \\
0.5 & 19.5 & 68.5 & 62.2 & 72.2 & 57.3 & 66.3 & 50.5 \\
0.7 & 15.9 & 69.0 & 62.8 & 73.5 & 57.4 & 66.9 & 52.1 \\
0.9 & 10.0 & 69.2 & 62.2 & 74.0 & 56.7 & 68.9 & 52.1 \\
\bottomrule
\end{tabular}
\end{table}

\paragraph{Q7. The estimated scale is too large on average. Can it be improved, and should VGG-T-Det refine it?}
VGGT-SLAM reconstructs a scene up to an unknown scale. The pipeline estimates the metric scale from the objects:
for every well-seen instance of a class with a size prior, the ratio of the prior footprint and height to the
measured ones is taken, and the scale is the 35th percentile of these ratios (a low percentile, because partly
seen objects measure too small and would otherwise inflate the scale).
\emph{Known so far:} against the true scale (from the ground-truth camera poses), the estimate is too large in
the median by 4\% and in the mean by 11\% (Table~\ref{tab:sam3-q7}). The error is skewed: a tenth of the scenes
is more than a third too large, a tenth more than 13\% too small. Only 48\% of the scenes are within 10\% and
74\% within 20\%. The scale error is the largest geometric error of the list: in the ablation of
Section~\ref{sec:sam3}, ground-truth box sizes are the only geometric correction that helps the language model,
and in the query check for the detector the estimated scale costs 9 to 14 points of box coverage compared with
the true scale.

\begin{table}[t]
\centering
\caption{Estimated scale divided by true scale, per scene (question names + ``room door'').}
\label{tab:sam3-q7}
\begin{tabular}{lrrrrrr}
\toprule
Dataset & Scenes & Median & Mean & 10th--90th percentile & Within 10\% & Within 20\% \\
\midrule
ScanNet     &  88 & 1.06 & 1.14 & 0.89--1.33 & 49\% & 78\% \\
ARKitScenes & 150 & 1.05 & 1.10 & 0.86--1.35 & 43\% & 70\% \\
ScanNet++   &  49 & 1.03 & 1.07 & 0.83--1.30 & 65\% & 78\% \\
All         & 287 & 1.04 & 1.11 & 0.87--1.34 & 48\% & 74\% \\
\bottomrule
\end{tabular}
\end{table}

Ways to improve the estimate (the first three are tested below):
\begin{itemize}
\item \emph{Camera height.} The height of the camera above the floor is measured in the reconstruction and is
      close to a constant for hand-held video; it gives a second, object-free scale estimate. It can be combined
      with the object ratios or used to reject outliers.
\item \emph{More reliable objects.} Use only instances seen completely (mask not cut by the image border, many
      views), prefer classes with a narrow size distribution (door height, toilet, bed length) over classes with
      a wide one (table, cabinet), and weight the ratios by the spread of the class prior.
\item \emph{Calibrate the percentile.} The 35th percentile was set by hand; the bias in Table~\ref{tab:sam3-q7}
      suggests a lower one, and it can be fitted on the training scenes, per dataset if needed.
\item \emph{A learned metric prior.} A monocular metric-depth model on a few keyframes gives the scale without
      object priors.
\end{itemize}

\emph{Tested without training} (Table~\ref{tab:sam3-q7b}; constants fitted on 939 ScanNet training scenes, evaluated on
the 287 VSI-Bench scenes; the ground truth is used for the evaluation only).
The bias is easy to remove: taking the 30th instead of the 35th percentile of the object ratios centres the
estimate (median 1.01) and halves the share of scenes that are more than 30\% too large (14\% to 9\%). The spread
hardly changes: the typical error stays near 15\% and only half of the scenes are within 10\%. The other ideas
do not help. The camera height is not a usable constant: its true value is 0.8\,m in ScanNet and 1.3\,m in
ARKitScenes and ScanNet++, and it varies from 0.5 to 1.4\,m within ScanNet, so a scale from the camera height
alone is far worse, and mixing it with the object estimate is worse than the object estimate alone. Weighting
the classes by how reliable their size is, with the per-class bias removed, gives the same error as the plain
percentile. The limit is therefore the size priors themselves: object sizes vary too much within a class to fix
the scale to better than about 15\% from a handful of objects. A better scale needs another source of metric
information, a metric-depth model or the detector.

\begin{table}[t]
\centering
\caption{Scale estimators that use no ground truth. Error: mean absolute log ratio of estimated to true scale.}
\label{tab:sam3-q7b}
\begin{tabular}{lrrrrr}
\toprule
Estimator & Median ratio & Error & Within 10\% & Within 20\% & More than 30\% too large \\
\midrule
Object ratios, 35th percentile (current)          & 1.04 & 0.154 & 48\% & 74\% & 14\% \\
Object ratios, 30th percentile                    & 1.01 & 0.146 & 51\% & 76\% &  9\% \\
Object ratios, 25th percentile                    & 0.98 & 0.152 & 49\% & 76\% &  7\% \\
Object heights only, 35th percentile              & 1.06 & 0.202 & 48\% & 71\% & 18\% \\
Classes weighted by reliability, bias removed     & 0.95 & 0.155 & 45\% & 75\% &  8\% \\
Camera height only                                & 0.90 & 0.367 & 14\% & 40\% & 31\% \\
Object estimate and camera height, equal weights  & 0.96 & 0.216 & 28\% & 61\% & 20\% \\
\bottomrule
\end{tabular}
\end{table}

Should VGG-T-Det refine the scale? Today the detector cannot: it is given a ground-truth scale (the mean distance
of the scan points from the first camera) with its input, in training and at test time, so it never sees the
scale error. Two designs are possible. (i) The detector keeps receiving a scale, now the estimated one, and its
size regression has to absorb the error; this needs training with the estimated scale, or with the true scale
perturbed by the error distribution of Table~\ref{tab:sam3-q7}. (ii) The detector predicts one scale correction
per scene from all its queries together, which is applied to every box; a per-scene quantity is better
estimated from all objects at once than box by box. In both cases the question is whether the images and VGGT
tokens carry enough metric information beyond the object size priors the pipeline already uses.


\subsection{Answers so far to Q1, Q2 and Q5}
\label{sec:sam3-answers}

\paragraph{Q1: coverage of the detector's training data.}
Table~\ref{tab:sam3-a1} counts the annotated objects that fall into the 65 classes. A third (ScanNet) to more
than half (ScanNet++) of the annotated objects are outside them; these are mostly small clutter (book, paper,
box, bottle, socket, clothes, bag, light switch), spread over 3,126 distinct labels in ScanNet++. Only
ScanNet++ has boxes for all 65 classes. Seven classes exist in ScanNet++ alone: ceiling light (2,501 boxes),
power strip (241), mattress (138), cutting board (113), paper bag (110), coat rack (71) and shoe rack (49).
ARKitScenes has no boxes for 14 classes, among them counter, desk, monitor, closet and computer tower.
Whether SAM 3 finds the annotated objects in the training scenes is not measured yet.

\begin{table}[t]
\centering
\caption{Annotated objects of the three training datasets and the share that falls into the 65 classes.
Structure: walls, floors, ceilings and the generic ``object''.}
\label{tab:sam3-a1}
\begin{tabular}{lrrrrr}
\toprule
Dataset & Annotations & Structure & Objects & In the 65 classes & Classes with boxes \\
\midrule
ScanNet (EmbodiedScan, 1,201 scenes)     &  48,964 & 30.1\% & 34,220 & 66.4\% & 57 \\
ARKitScenes (EmbodiedScan, 4,498 scenes) & 102,623 & 16.3\% & 85,926 & 86.9\% & 51 \\
ScanNet++ (905 scenes)                   & 100,799 & 29.9\% & 70,672 & 42.4\% & 65 \\
\bottomrule
\end{tabular}
\end{table}

\paragraph{Q2: coverage of the Reason-190K questions.}
The objects a training question is about are the labels of its object list (Table~\ref{tab:sam3-a2};
room-size questions name none and are left out). Questions from ScanNet and ARKitScenes are covered almost
completely; the only notable gap is bicycle (287 ScanNet questions). Of the ScanNet++ questions, 28\% name an
object outside the 65 classes, most often light switch (2,337 questions), board (2,089), bottle, box, switch,
bag, slippers, book, blinds and jacket. That is about 18,400 of the 143,437 training questions (13\%).

\begin{table}[t]
\centering
\caption{Reason-190K questions whose named objects are all in the 65 classes. Exact: the class names and their
VSI-Bench synonyms. With renamings: also the dataset labels mapped to a class (ceiling lamp, shelf, mouse,
coat hanger, shoe, poster, office chair, couch).}
\label{tab:sam3-a2}
\begin{tabular}{lrrrr}
\toprule
Source & Questions & Exact & With renamings & Object mentions covered \\
\midrule
ScanNet     &  49,569 & 98.7\% & 98.9\% & 99.5\% \\
ARKitScenes &  28,485 & 99.1\% & 99.6\% & 99.8\% \\
ScanNet++   &  65,383 & 64.4\% & 71.8\% & 87.0\% \\
All         & 143,437 & 83.1\% & 86.7\% & 93.6\% \\
\bottomrule
\end{tabular}
\end{table}

\paragraph{Q5, wording: the proposed rewordings mostly do not help.}
Each rewording of Table~\ref{tab:sam3-q5-few} was run as a single SAM 3 prompt on the scenes whose questions
name the object and substituted for the benchmark name (Table~\ref{tab:sam3-a5}). Only one is a clear gain:
prompting ``radiator'' also where the question says ``heater'' raises the boxes from 6 to 21 for 46 objects and
the objects found from 11\% to 39\%. ``Ceiling lamp'', ``desktop computer'', ``PC case'' and ``extension cord with sockets'' find
fewer objects than the benchmark name. ``Kitchen counter'' and ``bag'' find nearly all objects (85\% and 88\%)
but add many extra boxes (14 for 13 counters, 22 for 43 backpacks), so they exchange one count error for the
other. ``Ventilation fan'' changes nothing. The benchmark names are thus mostly as good as or better than the alternatives, and the low recall of
ceiling light, computer tower, fan and power strip is not a wording problem.

\begin{table}[t]
\centering
\caption{Rewordings tested (box score $\geq 0.7$). Found: objects with a box of that name within 0.75\,m.
Exact: scenes with the right number of boxes.}
\label{tab:sam3-a5}
\begin{tabular}{llrrrrrr}
\toprule
Object (ground truth) & Prompt & Boxes & Extra & Missing & Found & Exact & Verdict \\
\midrule
heater / radiator (46) & benchmark name            &  6 &  0 & 40 & 11\% & 18\% & \\
                       & ``radiator''              & 21 &  1 & 26 & 39\% & 44\% & better \\
                       & ``wall radiator''         & 15 &  1 & 32 & 28\% & 35\% & better \\
ceiling light (116)    & benchmark name            & 53 &  1 & 64 & 39\% & 24\% & \\
                       & ``ceiling lamp''          & 21 &  0 & 95 & 16\% & 10\% & worse \\
computer tower (42)    & benchmark name            & 25 &  1 & 18 & 50\% & 14\% & \\
                       & ``desktop computer''      & 37 & 10 & 15 & 38\% & 14\% & worse \\
                       & ``PC case''               & 13 &  0 & 29 & 26\% &  0\% & worse \\
fan / exhaust fan (14) & benchmark name            &  5 &  1 & 10 & 21\% & 30\% & \\
                       & ``ventilation fan''       &  6 &  1 &  9 & 29\% & 30\% & no change \\
power strip (10)       & benchmark name            &  6 &  1 &  5 & 40\% & 29\% & \\
                       & ``extension cord with sockets'' & 0 & 0 & 10 & 0\% & 0\% & worse \\
counter (13)           & benchmark name            & 12 &  6 &  7 & 31\% & 18\% & \\
                       & ``kitchen counter''       & 27 & 14 &  0 & 85\% & 36\% & more found, more extra \\
backpack (43)          & benchmark name            & 29 &  0 & 14 & 67\% & 52\% & \\
                       & ``bag''                   & 64 & 22 &  1 & 88\% & 55\% & more found, more extra \\
                       & ``backpack'' and ``bag''  & 68 & 25 &  0 & 93\% & 52\% & more found, more extra \\
\bottomrule
\end{tabular}
\end{table}

\paragraph{Q5, decision.}
``Radiator'' is adopted: every scene whose questions say ``heater'' or ``radiator'' is prompted with ``radiator'',
and the box keeps the name the question uses. Over all 287 scenes this lowers the missing boxes of the list from
641 to 627 and leaves the extra boxes unchanged (414 to 415); the share of (scene, name) pairs with the exact
count rises from 54.2\% to 54.8\%. Radiators remain a weak class: 26 of 46 are still missing.
``Kitchen counter'' and ``bag'' are on hold: they find nearly all objects but add many extra boxes, so they are
useful only once a later stage removes extra boxes. The other rewordings are rejected. The prompt set is
therefore the names of the scene's questions with two substitutions, ``room door'' for door and ``radiator''
for heater.

\paragraph{Q5, extra boxes: half of them are on nothing that is annotated.}
Of 2,967 boxes with score $\geq 0.7$ in all scenes, 2,368 match a ground-truth object of their name and 599 do
not. Of these 599, 299 (50\%) lie where nothing is annotated, 163 (27\%) duplicate an object that already has a
box, and 137 (23\%) lie on an object of another class (Table~\ref{tab:sam3-a5x}). This corrects two guesses of
Table~\ref{tab:sam3-q5-many}. Duplicates are not the main cause for table, chair and sofa, so better merging
would remove about a quarter of the extra boxes at most; only for bed do duplicates lead. And tv boxes are
rarely on monitors (2 cases). Chair and sofa are confused in both directions, and pillow and cushion boxes often
land on the sofa or bed they lie on. Whether a box on nothing annotated is a real object without annotation or a
false detection cannot be told from the counts.

\begin{table}[t]
\centering
\caption{What the boxes of a name are (all 287 scenes, score $\geq 0.7$).}
\label{tab:sam3-a5x}
\begin{tabular}{lrrrrrl}
\toprule
Name & Boxes & Matched & Duplicate & Other class & Nothing annotated & Main other class \\
\midrule
table   & 454 & 317 & 48 & 25 & 64 & chair 11, bed 5 \\
chair   & 864 & 739 & 23 & 27 & 75 & sofa 17, table 6 \\
sofa    & 144 & 106 & 11 & 10 & 17 & chair 8 \\
door    & 163 & 131 &  6 &  2 & 24 & -- \\
bed     & 106 &  77 & 14 &  2 & 13 & -- \\
pillow  &  93 &  70 &  9 & 13 &  1 & sofa 6, bed 4 \\
tv      &  96 &  78 &  8 &  5 &  5 & monitor 2 \\
cushion &  12 &   5 &  1 &  6 &  0 & sofa 3, chair 3 \\
\midrule
All names & 2,967 & 2,368 & 163 & 137 & 299 & \\
\bottomrule
\end{tabular}
\end{table}

\subsection{True scale against estimated scale}
\label{sec:sam3-scale}

The SAM 3 + VGGT-SLAM list estimates its metric scale from object sizes (Q7). To measure what that estimate costs, the
same list (prompts from the scene's questions, ``room door'', ``radiator'') was scored once as it is and once with
its scale replaced by the true one, obtained by fitting the reconstructed cameras to the ground-truth cameras.
Nothing else changes: the same boxes, only rescaled about the first camera. The estimated scale is a median 10\%
off (49\% of the scenes within 10\%, 87\% within 30\%).

\paragraph{3D boxes.} Table~\ref{tab:sam3-scale-3d}: predicted boxes are matched to ground-truth boxes of the same
label (287 VSI-Bench scenes, 3,194 objects); precision counts boxes with score $\geq 0.7$. The true scale improves
every loose-match measure by 5 to 9 points, on all three datasets (Table~\ref{tab:sam3-scale-ds}). Tight matches do
not improve (mAP@0.50 stays near 4): the scale only fixes the distance of every box from the camera, not the boxes'
shape and placement errors.

\begin{table}[t]
\centering
\caption{SAM 3 + VGGT-SLAM boxes with the estimated and the true scale (287 VSI-Bench scenes, 3,194 objects).}
\label{tab:sam3-scale-3d}
\begin{tabular}{lrrr}
\toprule
 & Estimated scale & True scale & Change \\
\midrule
Recall, IoU $\geq 0.25$             & 30.5 & 38.2 & $+7.7$ \\
Precision, IoU $\geq 0.25$          & 32.5 & 40.3 & $+7.8$ \\
mAP@0.25                            & 17.2 & 22.2 & $+5.0$ \\
Recall, IoU $\geq 0.5$              &  8.1 & 11.2 & $+3.1$ \\
mAP@0.50                            &  4.4 &  4.0 & $-0.4$ \\
Recall, centre within 0.5\,m        & 67.3 & 75.9 & $+8.6$ \\
Precision, centre within 0.5\,m     & 65.8 & 73.9 & $+8.1$ \\
\bottomrule
\end{tabular}
\end{table}

\begin{table}[t]
\centering
\caption{Recall at IoU $\geq 0.25$ by dataset.}
\label{tab:sam3-scale-ds}
\begin{tabular}{lrrr}
\toprule
Dataset & Objects & Estimated scale & True scale \\
\midrule
ScanNet     & 1,224 & 32.0 & 41.9 \\
ARKitScenes &   854 & 34.8 & 43.0 \\
ScanNet++   & 1,116 & 25.4 & 30.4 \\
\bottomrule
\end{tabular}
\end{table}

\paragraph{VSI-Bench.} Table~\ref{tab:sam3-scale-vsi} gives the score of the default language model (mix-v2-1) with
the two lists. The true scale helps the two metric tasks a little (absolute distance $+1.0$, size $+1.3$) and leaves
the overall score unchanged; counting even drops by 1.1, within the noise of a single run. The scale error therefore
costs far less than the errors in the list itself: which objects are listed and how loose their boxes are.

\begin{table}[t]
\centering
\caption{VSI-Bench (all 5,130 questions) with the SAM 3 + VGGT-SLAM list at the estimated and the true scale.}
\label{tab:sam3-scale-vsi}
\begin{tabular}{lrrrrr}
\toprule
SAM 3 list & Full & Debiased & Abs.\ distance & Size & Counting \\
\midrule
Estimated scale & 69.0 & 62.9 & 57.5 & 73.4 & 73.5 \\
True scale      & 69.1 & 63.0 & 58.5 & 74.7 & 72.4 \\
\bottomrule
\end{tabular}
\end{table}

\paragraph{Which results use the true scale.} Every VGG-T-Det variant (the released 18-class model, cls-31, cls-63,
cls-65 and IndoorDet-65-31k) multiplies VGGT's normalized point cloud by the mean distance of the ground-truth scan
from the first camera and uses the first camera's ground-truth pose, in training and at test time; the proposals of
IndoorDet-65-31k are placed with the true scale as well. All object lists built from a detector, and all language-model
training data (ground-truth and detector boxes), are therefore metric with the true scale. Only the SAM 3 + VGGT-SLAM
lists use an estimated scale, so they are the only video-only lists, and the language model sees a scale error only
at test time with them.

Table~\ref{tab:sam3-scale-use} lists which runs use the true scale.

\begin{table}[t]
\centering
\caption{Runs with and without the true scale. ``True scale'' = metric scale from the ground-truth scan (detectors) or
from fitting the reconstructed cameras to the ground-truth cameras (SAM 3 lists); ``estimated'' = from object size
priors, with up direction and walls from the reconstruction and the first camera as origin (no ground truth).}
\label{tab:sam3-scale-use}
\small
\begin{tabular}{llll}
\toprule
Run & Stage & Scale & Used for \\
\midrule
\multicolumn{4}{l}{\emph{Without the true scale (video only)}} \\
SAM 3 + VGGT-SLAM lists, all prompt variants (fixed 65 names, question & test & estimated & VSI-Bench, e.g.\ 69.0 / 62.9 \\
\quad names, ``room door'', ``radiator'', route phrases, thresholds 0.1--0.9) & & & \\
VGGT-SLAM lists with oracle 2D boxes (\texttt{slam\_manh}, \texttt{slamall\_manh}) & test & estimated & VSI-Bench \\
SAM 3 list ablation: predicted centre and size cells & test & estimated & VSI-Bench (Section~\ref{sec:sam3}) \\
SAM 3 list, estimated-scale arm of Table~\ref{tab:sam3-scale-3d} & test & estimated & 3D evaluation, VSI-Bench \\
Video-only model & train, test & no boxes & VSI-Bench 69.4 / 63.7 \\
\midrule
\multicolumn{4}{l}{\emph{With the true scale}} \\
VGG-T-Det: released 18-class, ES-56, cls-31, VSI-40, cls-63, cls-65 & train, test & true & detection; VSI-Bench lists \\
IndoorDet-65-31k (test2-hybrid), and its SAM 3 proposals & train, test & true & detection; VSI-Bench lists \\
Lists from a detector (cls-31; IndoorDet-65-31k refine and & test & true & VSI-Bench, e.g.\ 69.9 / 62.7, 69.8 / 63.2 \\
\quad threshold lists) & & & \\
VGGT-SLAM lists \texttt{\_truescale}; SAM 3 list, true-scale arm & test & true & VSI-Bench, 3D evaluation \\
Language-model SFT data (mix-v2-0, mix-v2-1, Exp1): ground-truth & train & true & every SFT run \\
\quad and detector boxes & & & \\
\bottomrule
\end{tabular}
\end{table}
